\documentclass[11pt,oneside]{article}

\usepackage[T1]{fontenc}
\usepackage{lmodern}
\usepackage{microtype}
\usepackage[margin=1in,headheight=14pt]{geometry}
\usepackage{amsmath,amssymb,amsthm,mathtools}
\usepackage{booktabs,longtable,tabularx,array}
\usepackage{enumitem}
\usepackage{xcolor}
\usepackage{xurl}
\usepackage{hyperref}
\usepackage{cleveref}
\usepackage{listings}
\usepackage{fancyhdr}
\usepackage{titlesec}
\usepackage{needspace}

\definecolor{provedgreen}{RGB}{16,102,70}
\definecolor{externalblue}{RGB}{25,87,143}
\definecolor{targetviolet}{RGB}{106,55,145}
\definecolor{warningred}{RGB}{145,45,45}
\definecolor{codegray}{RGB}{247,247,247}
\definecolor{rulegray}{RGB}{100,107,115}

\hypersetup{
  colorlinks=true,
  linkcolor=externalblue,
  citecolor=provedgreen,
  urlcolor=externalblue,
  pdftitle={Exponent-Lattice Moment Masks for the Alaoglu--Erdos Conjecture},
  pdfauthor={OpenAI GPT-5.6 Pro},
  pdfsubject={Toric reduction, moment masks, stable structural support, shared-pole Newton faces, and reduced denominators}
}
\urlstyle{same}

\pagestyle{fancy}
\fancyhf{}
\lhead{Alaoglu--Erd\H{o}s progress report}
\rhead{Exponent-lattice moment masks}
\cfoot{\thepage}
\renewcommand{\headrulewidth}{0.4pt}

\titleformat{\section}{\Large\bfseries}{\thesection}{0.7em}{}
\titleformat{\subsection}{\large\bfseries}{\thesubsection}{0.7em}{}
\titleformat{\subsubsection}{\normalsize\bfseries}{\thesubsubsection}{0.7em}{}

\setlength{\parindent}{0pt}
\setlength{\parskip}{0.55em}
\setlength{\emergencystretch}{3.5em}
\setlist[itemize]{leftmargin=2em,itemsep=0.25em,topsep=0.35em}
\setlist[enumerate]{leftmargin=2.3em,itemsep=0.25em,topsep=0.35em}

\newtheorem{theorem}{Theorem}[section]
\newtheorem{proposition}[theorem]{Proposition}
\newtheorem{lemma}[theorem]{Lemma}
\newtheorem{corollary}[theorem]{Corollary}
\newtheorem{conjecture}[theorem]{Conjecture}
\newtheorem{criterion}[theorem]{Criterion}
\newtheorem{question}[theorem]{Question}
\theoremstyle{definition}
\newtheorem{definition}[theorem]{Definition}
\newtheorem{example}[theorem]{Example}
\newtheorem{experiment}[theorem]{Computational experiment}
\newtheorem{target}[theorem]{Research target}
\newtheorem{warning}[theorem]{Warning}
\theoremstyle{remark}
\newtheorem{remark}[theorem]{Remark}

\newcommand{\N}{\mathbb N}
\newcommand{\Nzero}{\mathbb N_0}
\newcommand{\Z}{\mathbb Z}
\newcommand{\Q}{\mathbb Q}
\newcommand{\R}{\mathbb R}
\newcommand{\C}{\mathbb C}
\newcommand{\vp}{v_p}
\newcommand{\ord}{\operatorname{ord}}
\newcommand{\supp}{\operatorname{supp}}
\newcommand{\den}{\operatorname{den}}
\newcommand{\num}{\operatorname{num}}
\newcommand{\lcm}{\operatorname{lcm}}
\newcommand{\cont}{\operatorname{cont}}
\newcommand{\rad}{\operatorname{rad}}
\newcommand{\Span}{\operatorname{span}}
\newcommand{\rank}{\operatorname{rank}}
\newcommand{\AEname}{Alaoglu--Erd\H{o}s}
\newcommand{\thetaAE}{\vartheta}
\newcommand{\fall}[2]{\left(#1\right)_{\!\underline{#2}}}
\newcommand{\normone}[1]{\left\lVert #1\right\rVert_1}
\newcommand{\statusproved}{\textcolor{provedgreen}{\textbf{[proved here]}}}
\newcommand{\statusexternal}{\textcolor{externalblue}{\textbf{[external theorem]}}}
\newcommand{\statustarget}{\textcolor{targetviolet}{\textbf{[open target]}}}
\newcommand{\statuscomp}{\textcolor{targetviolet}{\textbf{[exact computation]}}}
\newcommand{\statuswarning}{\textcolor{warningred}{\textbf{[audit warning]}}}

\newenvironment{statusbox}[1]
{\begin{center}\begin{minipage}{0.94\textwidth}\raggedright\hrule\vspace{0.45em}
\textbf{Status:} #1\par\vspace{0.35em}}
{\vspace{0.35em}\hrule\end{minipage}\end{center}}

\newenvironment{resultbox}
{\begin{center}\begin{minipage}{0.94\textwidth}\hrule\vspace{0.45em}}
{\vspace{0.35em}\hrule\end{minipage}\end{center}}

\lstdefinestyle{pythonstyle}{
  language=Python,
  basicstyle=\ttfamily\scriptsize,
  backgroundcolor=\color{codegray},
  frame=single,
  rulecolor=\color{rulegray},
  numbers=left,
  numberstyle=\tiny\color{rulegray},
  numbersep=7pt,
  showstringspaces=false,
  breaklines=true,
  columns=fullflexible,
  keepspaces=true,
  tabsize=4
}

\begin{document}
\pagenumbering{gobble}

\begin{titlepage}
\centering
\vspace*{0.7cm}
{\Huge\bfseries Exponent-Lattice Moment Masks\par}
\vspace{0.2cm}
{\Huge\bfseries for the \AEname\ Conjecture\par}
\vspace{0.65cm}
{\LARGE Toric Collapse, Stable Structural Support,\par}
\vspace{0.1cm}
{\LARGE Shared-Pole Newton Faces, and the Reduced-Denominator Endgame\par}
\vspace{1.0cm}
{\large A nonduplicative progress report continuing the supplied unified manuscript\par}
\vspace{0.35cm}
{\large 22 August 2026\par}
\vfill
\begin{minipage}{0.91\textwidth}
\small
\textbf{Bottom line.}
No unconditional proof is claimed.  The new progress is a self-contained arithmetic framework
for a class of auxiliary constructions not developed in the supplied report: integer
``moment masks'' in exponent-lattice near-unit coordinates.  The framework produces five
exact conclusions.

First, a collection of arbitrarily many continued-fraction near-relations has only two
formal degrees of freedom: after quotienting by its toric lattice ideal, every mask is a
Laurent polynomial in two adjacent near-units.  Continued-fraction recurrences themselves
therefore produce identically zero masks at both the input and output, not small nonzero
integers.

Second, every one-direction Prouhet or Thue--Morse cancellation of order $r$ clears to an
integer divisible by the $r$th power of the corresponding output gap.  High contact on one
ray is thus not merely expensive; it algebraically factorizes into a large integer gap.

Third, a bivariate mask of contact order $r$ clears into the ideal generated by the two
output gaps to the $r$th power.  An exact primewise Newton functional gives the guaranteed
integer divisor and a product-formula contradiction criterion.

Fourth, the private-prime leading-face law is exact but not asymptotically decisive.
For adjacent continued-fraction near-units, every structural prime has an eventually fixed
sign: the reduced numerator and denominator supports stabilize, and both denominators are
supported on the same nonempty set of primes.  Their pole depths grow linearly and satisfy
exact determinant-one transport identities.

Fifth, at each shared denominator prime the cleared mask is governed by an explicit
two-dimensional Newton polygon.  A unique minimizing monomial survives exactly; cancellation
requires a tied initial form to vanish, and deep compression requires this to repeat through
successive $p$-adic layers at every stable denominator prime.  Together with effective
two-logarithm separation, this proves that visible denominator clearing cannot become subunit
and isolates the only surviving mechanism: near-total \emph{reduced-denominator cancellation}
through simultaneous, prime-specific shared Newton faces.  The report ends with an exact
sufficient criterion, a finite shared-face capture target, and a Lean-oriented dependency
graph.
\end{minipage}
\vfill
{\small Prepared by OpenAI GPT-5.6 Pro.}
\end{titlepage}

\pagenumbering{roman}
\tableofcontents
\clearpage
\pagenumbering{arabic}

\section{Scope, nonduplication boundary, and evidence convention}
\label{sec:scope}

The supplied unified report is exceptionally broad.  It already develops the classical
Four Exponentials reduction; the exact conditional solution monoid; valuation lattices,
Kummer theory and local common exponents; generalized Vandermonde and Loewner determinants;
finite differences, Newton interpolation and Smith profiles; Sturmian and carry models;
Mahler and special-function transfers; canonical products and Stieltjes moments; rational
point counting; support theorems; and multiple exact failure analyses.  Repeating any of
those surveys would add volume rather than information.

This supplement therefore imposes the following boundary.

\begin{itemize}
\item The exact conditional classification, finite exclusion ranges, Four Exponentials
  barrier, and all named results of the unified report are treated as inherited context and
  are not reproved.
\item The new variables are \emph{near-units attached to exponent-lattice relations}, not
  the original smooth nodes.  The auxiliary objects are integer masks in those variables.
\item Determinants appear only as the $2\times2$ lattice determinant controlling two
  independent near-relations.  No new Vandermonde or interpolation determinant is proposed.
\item The main arithmetic object is a single cleared rational value and its reduced
  denominator.  This is intentionally scalar.
\item Every assertion is marked as proved here, an external input, an exact computation, or
  an open target.  The conjecture remains open at the end.
\end{itemize}

The public-claims audit in \cref{sec:public-audit} confirms that several recent machine-assisted
mathematical results mentioned in the prompt have public records, but it found no public
technical statement of the alleged new \AEname\ step.  Nothing from that unavailable claim is
used below.

\begin{statusbox}{\statusproved{} for the mathematical statements proved in this report;
\statusexternal{} for explicitly cited inputs; \statuscomp{} for the small exact symbolic and
continued-fraction computations; \statustarget{} for the proposed endgame.}
\end{statusbox}

\section{Inherited setup and the relation lattice}
\label{sec:setup}

Write
\[
  \thetaAE=\frac{\log 3}{\log 2}.
\]
Suppose, for contradiction, that a nonintegral real number $x$ satisfies
\begin{equation}
  M:=2^x\in\N,
  \qquad
  N:=3^x\in\N.
  \label{eq:MN}
\end{equation}
The elementary interval argument gives $x>1$: for $x<0$ one has $0<2^x<1$, and for
$0<x<1$ one has $1<2^x<2$.  The supplied report proves much more, including transcendence of
$x$ and an exact primitive-generator theorem, but only the following elementary fact is
needed here.

\begin{lemma}[Output multiplicative independence]
\label{lem:MN-independent}
The integers $M$ and $N$ are multiplicatively independent.
\end{lemma}

\begin{proof}
If $M^aN^b=1$ for $a,b\in\Z$, then taking real logarithms and using
\eqref{eq:MN} gives
\[
  x(a\log2+b\log3)=0.
\]
Since $x\ne0$, unique factorization applied to $2^a3^b=1$ gives $a=b=0$.
\end{proof}

For an exponent vector $\nu=(a,b)\in\Z^2$, define its input and output characters
\begin{equation}
 \rho_\nu=2^a3^b,
 \qquad
 \sigma_\nu=M^aN^b=\rho_\nu^x,
 \qquad
 \delta_\nu=a\log2+b\log3.
 \label{eq:characters}
\end{equation}
Thus $\rho_\nu=e^{\delta_\nu}$ and $\sigma_\nu=e^{x\delta_\nu}$.  A good rational
approximation to $\thetaAE$ supplies a nonzero $\nu$ for which $\delta_\nu$ is small; then
both $\rho_\nu$ and the rational number $\sigma_\nu$ are close to $1$.

If $p/q$ is a convergent to $\thetaAE$, the natural vector is
\begin{equation}
  \nu=(-p,q),
  \qquad
  \delta_\nu=q\log3-p\log2.
  \label{eq:cf-vector}
\end{equation}
The central question of this report is whether several such rational output near-units can
be combined by an integer polynomial mask to produce a nonzero rational number that is too
small for its reduced denominator.

\subsection{Exact reduced-height growth}
\label{subsec:height}

The denominator cost of a rational output near-unit is controlled by the valuation geometry
of $(M,N)$, independently of continued fractions.

Let $\mathcal P$ be the finite set of primes dividing $MN$, and put
\[
  \mathbf m=(v_p(M))_{p\in\mathcal P},
  \qquad
  \mathbf n=(v_p(N))_{p\in\mathcal P}.
\]
For $\nu=(a,b)$ define the weighted valuation norm
\begin{equation}
  L_{M,N}(\nu)
   =\sum_{p\in\mathcal P}|a v_p(M)+b v_p(N)|\log p.
  \label{eq:valuation-norm}
\end{equation}

\begin{lemma}[Valuation norm]
\label{lem:valuation-norm}
The function $L_{M,N}$ extends to a norm on $\R^2$.  Consequently there is a constant
$c_{M,N}>0$ such that
\begin{equation}
  L_{M,N}(a,b)\ge c_{M,N}(|a|+|b|)
  \qquad(a,b\in\R).
  \label{eq:norm-equivalence}
\end{equation}
\end{lemma}

\begin{proof}
It is a seminorm by construction.  If it vanishes at $(a,b)$, then
$a\mathbf m+b\mathbf n=0$.  For rational $(a,b)$ this would give a multiplicative relation
between $M$ and $N$; hence the two valuation vectors are linearly independent over $\Q$.
Because they have integral coordinates, they are also linearly independent over $\R$.
Thus the seminorm has trivial kernel.  Norm equivalence in the finite-dimensional space
$\R^2$ gives \eqref{eq:norm-equivalence}.
\end{proof}

Write the rational number $\sigma_\nu$ in lowest terms as
\[
  \sigma_\nu=\frac{R_\nu}{S_\nu},
  \qquad R_\nu,S_\nu\in\N,
  \qquad \gcd(R_\nu,S_\nu)=1.
\]

\begin{proposition}[Exact numerator--denominator law]
\label{prop:height-law}
For every $\nu\in\Z^2$,
\begin{align}
 \log R_\nu+\log S_\nu&=L_{M,N}(\nu),
 \label{eq:height-sum}\\
 \log R_\nu-\log S_\nu&=x\delta_\nu.
 \label{eq:height-difference}
\end{align}
In particular,
\begin{equation}
  h(\sigma_\nu):=\log\max(R_\nu,S_\nu)
  =\frac{L_{M,N}(\nu)+|x\delta_\nu|}{2},
  \label{eq:height-exact}
\end{equation}
and if $|x\delta_\nu|\le1$, then
\begin{equation}
 \min(\log R_\nu,\log S_\nu)
 \ge \frac{c_{M,N}(|a|+|b|)-1}{2}.
 \label{eq:both-sides-large}
\end{equation}
\end{proposition}

\begin{proof}
The exponent of a prime $p$ in $\sigma_\nu$ is
$e_p=a v_p(M)+b v_p(N)$.  Reduction separates the positive and negative parts:
\[
 \log R_\nu=\sum_p\max(e_p,0)\log p,
 \qquad
 \log S_\nu=\sum_p\max(-e_p,0)\log p.
\]
Their sum is \eqref{eq:height-sum}; their difference is
$\log\sigma_\nu=x\delta_\nu$, giving \eqref{eq:height-difference}.  The final two formulas
follow by adding and subtracting.
\end{proof}

The significance is quantitative.  A nontrivial rational output near-unit can be
archimedeanly close to $1$, but its reduced numerator and denominator are both exponentially
large in the exponent-vector length.  Any auxiliary mask must defeat that height, not merely
make a Taylor series start late.

\subsection{A two-direction uncertainty identity}
\label{subsec:uncertainty}

Let
\[
 \nu=(a,b),\qquad \omega=(c,d),\qquad
 D=\det\begin{pmatrix}a&b\\c&d\end{pmatrix}=ad-bc.
\]

\begin{proposition}[Exact relation-lattice uncertainty]
\label{prop:uncertainty}
One has
\begin{equation}
 d\delta_\nu-b\delta_\omega=D\log2,
 \qquad
 a\delta_\omega-c\delta_\nu=D\log3.
 \label{eq:uncertainty-identities}
\end{equation}
Consequently, if $D\ne0$ and
$H=\normone{\nu}+\normone{\omega}$, then
\begin{equation}
 \max(|\delta_\nu|,|\delta_\omega|)
 \ge \frac{|D|\max(\log2,\log3)}{H}.
 \label{eq:uncertainty-bound}
\end{equation}
\end{proposition}

\begin{proof}
Substitute $\delta_\nu=a\log2+b\log3$ and
$\delta_\omega=c\log2+d\log3$.  The cross terms cancel and give
\eqref{eq:uncertainty-identities}.  Each right-hand side is bounded above in absolute value
by $H\max(|\delta_\nu|,|\delta_\omega|)$.
\end{proof}

For adjacent convergents $p_n/q_n,p_{n+1}/q_{n+1}$ to $\thetaAE$, the vectors
$\nu_n=(-p_n,q_n)$ and $\nu_{n+1}=(-p_{n+1},q_{n+1})$ have determinant $\pm1$.
They are therefore the optimally conditioned integral basis for two-direction near-unit
constructions.  Formula \eqref{eq:uncertainty-bound} says that even this basis cannot make
both logarithmic errors smaller than the reciprocal exponent scale.

\begin{statusbox}{\statusproved{} for \cref{lem:MN-independent,lem:valuation-norm,prop:height-law,prop:uncertainty}.
These are elementary and independent of the deep transcendence inputs in the unified report.}
\end{statusbox}

\section{Toric reduction: many near-relations have only two formal degrees of freedom}
\label{sec:toric}

A natural reaction to the two-direction uncertainty bound is to use many convergents and
Prouhet-type cancellations among them.  The first new structural result shows exactly what
additional variables do and do not buy.

Fix exponent vectors $\nu_i=(a_i,b_i)\in\Z^2$ for $1\le i\le k$.  Define the formal monomial
map
\begin{equation}
 \Psi_\nu:\Z[T_1^{\pm1},\ldots,T_k^{\pm1}]
 \longrightarrow \Z[U^{\pm1},V^{\pm1}],
 \qquad
 T_i\longmapsto U^{a_i}V^{b_i}.
 \label{eq:monomial-map}
\end{equation}
For $\alpha=(\alpha_1,\ldots,\alpha_k)\in\Z^k$ write
$T^\alpha=\prod_iT_i^{\alpha_i}$ and
\[
 \pi_\nu(\alpha)=\sum_i\alpha_i\nu_i\in\Z^2.
\]

\begin{definition}[Toric lattice ideal]
\label{def:toric-ideal}
The lattice ideal of the relation list $\nu$ is
\begin{equation}
 I_\nu=
 \left\langle T^\alpha-T^\beta:
 \pi_\nu(\alpha)=\pi_\nu(\beta)\right\rangle
 \subset\Z[T_1^{\pm1},\ldots,T_k^{\pm1}].
 \label{eq:toric-ideal}
\end{equation}
\end{definition}

\begin{theorem}[Universal toric kernel]
\label{thm:toric-kernel}
The kernel of $\Psi_\nu$ is exactly $I_\nu$.
\end{theorem}

\begin{proof}
Every displayed binomial has the same aggregate exponent on its two terms, so
$I_\nu\subseteq\ker\Psi_\nu$.  Conversely, let
$C=\sum_\alpha c_\alpha T^\alpha$ lie in the kernel.  Group its terms by the fibers of
$\pi_\nu$.  Laurent monomials $U^rV^s$ are a free $\Z$-basis of the target ring, so in each
fiber the sum of the coefficients is zero.  Choose one reference exponent $\alpha_0$ in a
fiber.  Then
\[
 \sum_{\alpha\text{ in the fiber}}c_\alpha T^\alpha
 =\sum_{\alpha\ne\alpha_0}c_\alpha
   (T^\alpha-T^{\alpha_0}),
\]
because the coefficient of $T^{\alpha_0}$ is the negative sum of the others.  Every
difference belongs to $I_\nu$.
\end{proof}

\begin{corollary}[Universal identities transfer to the output]
\label{cor:toric-transfer}
If $C\in I_\nu$, then
\begin{equation}
 C(\rho_{\nu_1},\ldots,\rho_{\nu_k})=0,
 \qquad
 C(\sigma_{\nu_1},\ldots,\sigma_{\nu_k})=0.
 \label{eq:toric-both}
\end{equation}
\end{corollary}

\begin{proof}
Compose \eqref{eq:monomial-map} first with $(U,V)=(2,3)$ and then with $(U,V)=(M,N)$.
\end{proof}

\begin{warning}[The theorem is formal, not a classification of numerical additive identities]
\label{warn:numeric-kernel}
The numerical evaluation kernel at $(2,3)$ is larger than $I_\nu$, because $3-2-1=0$ and
similar additive identities exist.  The theorem classifies identities caused solely by
integer dependencies of exponent vectors.  Such identities are precisely the ones that
transfer automatically to every pair of bases.  Candidate-specific additive cancellation is
not ruled out; it is the difficult part.
\end{warning}

\subsection{Unimodular two-variable normal form}

\begin{theorem}[Two-variable normal form]
\label{thm:two-variable-normal-form}
Suppose two vectors $\nu_1,\nu_2$ form a unimodular basis of $\Z^2$.  For every
$i\ge3$ there are unique integers $r_i,s_i$ such that
$\nu_i=r_i\nu_1+s_i\nu_2$.  Modulo $I_\nu$, every mask has a unique Laurent normal form
\begin{equation}
 F(U,V)\in\Z[U^{\pm1},V^{\pm1}],
 \qquad
 T_1\mapsto U,
 \quad T_2\mapsto V,
 \quad T_i\mapsto U^{r_i}V^{s_i}.
 \label{eq:two-variable-normal-form}
\end{equation}
In particular, adding arbitrarily many convergent variables does not increase the formal
dimension beyond two.
\end{theorem}

\begin{proof}
Unimodularity makes $(\nu_1,\nu_2)$ a $\Z$-basis, giving the unique coordinates.  The
monomial map is then surjective onto a Laurent polynomial ring in the two basis monomials,
and \cref{thm:toric-kernel} identifies its kernel.
\end{proof}

\begin{example}[Continued-fraction recurrence masks are identically zero]
\label{ex:cf-toric-zero}
Continued-fraction vectors satisfy
\[
 \nu_{n+1}=a_{n+1}\nu_n+\nu_{n-1}.
\]
Therefore
\begin{equation}
 T_{n+1}-T_n^{a_{n+1}}T_{n-1}\in I_\nu.
 \label{eq:cf-toric-binomial}
\end{equation}
The corresponding input and output identities are exact.  Iterating the recurrence can
create impressive formal cancellation, but every such recurrence-only mask evaluates to
zero rather than to a useful nonzero small integer.
\end{example}

This is the first principal no-go result of the report.  A many-convergent construction must
be reduced modulo $I_\nu$ before its contact order is credited.  Otherwise torically trivial
zeros can be mistaken for high-order arithmetic cancellation.

\begin{statusbox}{\statusproved{} for the universal toric kernel, its transfer to both base
pairs, and the unimodular normal form.  The result is an elementary special case of the theory
of lattice and binomial ideals, included with a complete proof so that no black box is needed.}
\end{statusbox}

\section{Moment masks and contact order at the identity}
\label{sec:moment-masks}

After toric reduction, the right object is a Laurent polynomial in two near-unit variables.
Multiplying by a monomial changes neither vanishing order at $(1,1)$ nor the essential
arithmetic problem, so we may work with an ordinary polynomial
\[
 F(U,V)=\sum_{(i,j)\in\Lambda}c_{ij}U^iV^j\in\Z[U,V],
 \qquad c_{ij}\ne0.
\]

\begin{definition}[Contact order]
\label{def:contact-order}
The contact order of a nonzero mask $F$ at the identity is
\begin{equation}
 \ord_{(1,1)}F
 =\min\{r+s:[X^rY^s]F(1+X,1+Y)\ne0\}.
 \label{eq:contact-order}
\end{equation}
Equivalently, it is the largest integer $m$ such that
$F\in(U-1,V-1)^m$ in $\Z[U,V]$.
\end{definition}

If the contact order is $r$, evaluation at $U=e^{\varepsilon_1}$,
$V=e^{\varepsilon_2}$ begins at total degree $r$ in the two small logarithmic errors.  The
terminology ``moment mask'' comes from the following exact characterization.

\begin{proposition}[Two-dimensional Prouhet moment equations]
\label{prop:moment-equations}
Let
$F(U,V)=\sum_{\ell=1}^s c_\ell U^{a_\ell}V^{b_\ell}$ with distinct exponent points
$(a_\ell,b_\ell)$.  The following are equivalent for an integer $r\ge1$.
\begin{enumerate}
\item $\ord_{(1,1)}F\ge r$.
\item For all $i,j\ge0$ with $i+j<r$,
\begin{equation}
  \sum_{\ell=1}^s c_\ell
   \fall{a_\ell}{i}\fall{b_\ell}{j}=0.
  \label{eq:falling-moments}
\end{equation}
\item For all $i,j\ge0$ with $i+j<r$,
\begin{equation}
  \sum_{\ell=1}^s c_\ell a_\ell^i b_\ell^j=0.
  \label{eq:ordinary-moments}
\end{equation}
\end{enumerate}
\end{proposition}

\begin{proof}
The mixed derivative $\partial_U^i\partial_V^jF(1,1)$ is the left side of
\eqref{eq:falling-moments}.  Vanishing of all derivatives of total order below $r$ is
precisely membership in $(U-1,V-1)^r$.  Falling factorials and ordinary powers are related by
unitriangular Stirling transforms in each variable, so the two families of moment equations
are equivalent.
\end{proof}

Thus a contact-$r$ mask is a signed, weighted two-dimensional Prouhet--Tarry--Escott
configuration on the exponent lattice.  The next bounds quantify the unavoidable size of
such a configuration.

\begin{theorem}[Degree and support ceilings]
\label{thm:contact-ceilings}
Let $F\ne0$ have bidegree $(d,e)$ and $s$ nonzero monomials.  Then
\begin{equation}
 \ord_{(1,1)}F\le d+e,
 \qquad
 \ord_{(1,1)}F\le s-1.
 \label{eq:contact-ceilings}
\end{equation}
\end{theorem}

\begin{proof}
The polynomial $F(1+X,1+Y)$ has total degree at most $d+e$, giving the first bound.
For the second, write the exponent points as $(a_\ell,b_\ell)$.  Choose real numbers
$\alpha,\beta$ so that the frequencies
$\lambda_\ell=\alpha a_\ell+\beta b_\ell$ are distinct and so that the first nonzero
homogeneous part of $F(1+X,1+Y)$ does not vanish on the tangent direction
$(\alpha,\beta)$.  Then
\[
 g(t)=F(e^{\alpha t},e^{\beta t})
     =\sum_{\ell=1}^s c_\ell e^{\lambda_\ell t}
\]
has a zero at $0$ of exactly the same order as $F$ at $(1,1)$.  If that order were at least
$s$, then
$g^{(m)}(0)=\sum_\ell c_\ell\lambda_\ell^m=0$ for $0\le m<s$.  The Vandermonde matrix on the
distinct $\lambda_\ell$ is invertible, forcing every $c_\ell=0$, a contradiction.
\end{proof}

The support ceiling is a useful audit rule.  A claimed order-$r$ mask with at most $r$
monomials is either wrong or torically zero after normalization.

\subsection{A uniform local estimate}
\label{subsec:local-estimate}

Write
\begin{equation}
 \widetilde F(X,Y)=F(1+X,1+Y)=\sum_{i,j}\gamma_{ij}X^iY^j,
 \qquad
 \|\widetilde F\|_1=\sum_{i,j}|\gamma_{ij}|.
 \label{eq:shifted-mask}
\end{equation}

\begin{lemma}[Contact estimate]
\label{lem:contact-estimate}
Suppose $\ord_{(1,1)}F\ge r$ and
$|\varepsilon_1|,|\varepsilon_2|\le\tfrac14$.  Put
$\Delta=\max(|\varepsilon_1|,|\varepsilon_2|)$.  Then
\begin{equation}
 |F(e^{\varepsilon_1},e^{\varepsilon_2})|
 \le \|\widetilde F\|_1(2\Delta)^r.
 \label{eq:contact-estimate}
\end{equation}
\end{lemma}

\begin{proof}
For $|t|\le1/4$, one has $|e^t-1|\le2|t|$.  Every nonzero term in
\eqref{eq:shifted-mask} has $i+j\ge r$, and $2\Delta\le1$, so
$|e^{\varepsilon_1}-1|^i|e^{\varepsilon_2}-1|^j\le(2\Delta)^r$.
Sum the absolute values.
\end{proof}

The estimate is deliberately crude but universal.  Later we prove that this entire style of
estimate, when followed by full denominator clearing, cannot close the conjecture on
adjacent convergents.  Its failure identifies reduced-denominator cancellation as the only
remaining scalar-mask mechanism.

\section{One-direction Prouhet cancellation is an exact gap power}
\label{sec:one-direction}

Before using two variables, it is important to close the seemingly attractive one-variable
route completely.

\begin{theorem}[One-direction gap-power theorem]
\label{thm:univariate-gap-power}
Let $P(T)\in\Z[T]$ have degree $d$ and a zero of multiplicity at least $r$ at $T=1$.
For integers $R,S$ with $S\ne0$, define
\begin{equation}
 I_P(R,S)=S^dP(R/S)\in\Z.
 \label{eq:univariate-cleared}
\end{equation}
Then
\begin{equation}
 (R-S)^r\mid I_P(R,S).
 \label{eq:gap-power-divides}
\end{equation}
If $P(R/S)\ne0$, then
\begin{equation}
 |I_P(R,S)|\ge |R-S|^r.
 \label{eq:gap-power-floor}
\end{equation}
\end{theorem}

\begin{proof}
Because $T-1$ is monic, repeated integral polynomial division gives
$P(T)=(T-1)^rQ(T)$ with $Q\in\Z[T]$ and $\deg Q=d-r$.  Hence
\[
 S^dP(R/S)
 =(R-S)^r S^{d-r}Q(R/S),
\]
and the final factor is an integer.  The lower bound follows if the product is nonzero.
\end{proof}

This theorem covers every one-dimensional Prouhet partition, finite-difference mask, and
Thue--Morse sign pattern at once.  Coefficient cancellation cannot make the cleared integer
smaller than the corresponding output gap power.

\subsection{The Thue--Morse calibration}

For $r\ge1$, define
\begin{equation}
 P_r(T)=\prod_{j=0}^{r-1}(1-T^{2^j})
       =\sum_{n=0}^{2^r-1}(-1)^{s_2(n)}T^n,
 \label{eq:thue-morse-mask}
\end{equation}
where $s_2(n)$ is the sum of the binary digits of $n$.  It has degree $2^r-1$, exactly
$2^r$ coefficients of absolute value $1$, and a zero of exact multiplicity $r$ at $1$.
The corresponding Prouhet identity equalizes power sums of degrees $0,\ldots,r-1$ between
the even and odd Thue--Morse classes.

\begin{corollary}[Exact Thue--Morse factorization]
\label{cor:thue-morse-factorization}
For $d=2^r-1$,
\begin{equation}
 S^dP_r(R/S)
 =\prod_{j=0}^{r-1}\bigl(S^{2^j}-R^{2^j}\bigr).
 \label{eq:thue-morse-cleared}
\end{equation}
Every factor on the right is divisible by $S-R$, so the cleared value is divisible by
$(S-R)^r$.
\end{corollary}

\begin{proof}
The exponents $2^j$ sum to $2^r-1$, so denominator clearing distributes exactly across the
product in \eqref{eq:thue-morse-mask}.
\end{proof}

\begin{table}[ht]
\centering
\caption{Exact size data for the first Thue--Morse masks.}
\label{tab:thue-morse}
\begin{tabular}{@{}rrrr@{}}
\toprule
contact $r$ & degree $2^r-1$ & coefficient $\ell^1$ norm & multiplicity at $1$\\
\midrule
1&1&2&1\\
2&3&4&2\\
3&7&8&3\\
4&15&16&4\\
5&31&32&5\\
6&63&64&6\\
7&127&128&7\\
\bottomrule
\end{tabular}
\end{table}

The high contact is real, but the arithmetic output becomes a product of increasingly large
integer differences.  This is the opposite of a subunit certificate.

\begin{resultbox}
\textbf{Closed route.}  No mask supported on integer powers of a single rational output
near-unit can prove the conjecture by ``high Prouhet contact plus denominator clearing.''  If
the cleared value is nonzero, its exact gap-power divisor grows with the contact order.
\end{resultbox}

\section{Bivariate clearing: the gap ideal and its primewise Newton functional}
\label{sec:gap-ideal}

Two independent near-unit directions can distribute contact among two gaps and thereby evade
the individual gap power in \cref{thm:univariate-gap-power}.  The replacement is an exact
ideal-membership statement.

Let
\[
 z_1=\frac{R_1}{S_1},\qquad
 z_2=\frac{R_2}{S_2},
 \qquad
 G_1=R_1-S_1,\qquad G_2=R_2-S_2,
\]
where at first no coprimality assumption is needed.  Let $F\in\Z[U,V]$ have bidegree at most
$(d,e)$ and shifted expansion
\begin{equation}
 F(1+X,1+Y)=\sum_{\substack{0\le i\le d\\0\le j\le e}}\gamma_{ij}X^iY^j.
 \label{eq:shifted-coeffs}
\end{equation}
Define the fully cleared integer
\begin{equation}
 J_F(z_1,z_2)=S_1^dS_2^eF(z_1,z_2)\in\Z.
 \label{eq:cleared-bivariate}
\end{equation}

\begin{theorem}[Gap-ideal theorem]
\label{thm:gap-ideal}
If $\ord_{(1,1)}F\ge r$, then
\begin{equation}
 J_F(z_1,z_2)
 =\sum_{i+j\ge r}\gamma_{ij}
   G_1^iG_2^jS_1^{d-i}S_2^{e-j}
 \in (G_1,G_2)^r\subset\Z.
 \label{eq:gap-ideal-expansion}
\end{equation}
Consequently,
\begin{equation}
 \gcd(G_1,G_2)^r\mid J_F(z_1,z_2).
 \label{eq:gcd-gap-floor}
\end{equation}
\end{theorem}

\begin{proof}
Substitute $X=G_1/S_1$ and $Y=G_2/S_2$ in \eqref{eq:shifted-coeffs} and multiply by
$S_1^dS_2^e$.  Contact order removes all terms with $i+j<r$.  Every remaining monomial lies
in $(G_1,G_2)^r$, and every one is divisible by the $r$th power of their integer gcd.
\end{proof}

The gcd lower bound is often far from the full ideal information.  Prime by prime, the
support of the shifted mask determines a sharper guaranteed divisor.

\begin{definition}[Primewise Newton floor]
\label{def:newton-floor}
For a prime $p$, write
\[
 g_{k,p}=v_p(G_k),\qquad s_{k,p}=v_p(S_k)\quad(k=1,2).
\]
Define
\begin{equation}
 \tau_p(F;z_1,z_2)
 =\min_{\gamma_{ij}\ne0}
 \bigl(
 v_p(\gamma_{ij})+i g_{1,p}+j g_{2,p}
 +(d-i)s_{1,p}+(e-j)s_{2,p}
 \bigr).
 \label{eq:newton-floor}
\end{equation}
\end{definition}

\begin{proposition}[Exact guaranteed divisor]
\label{prop:guaranteed-divisor}
For every prime $p$,
\begin{equation}
 v_p(J_F(z_1,z_2))\ge\tau_p(F;z_1,z_2).
 \label{eq:vp-guaranteed}
\end{equation}
Thus the finite product
\begin{equation}
 \mathcal D_F(z_1,z_2)=\prod_p p^{\tau_p(F;z_1,z_2)}
 \label{eq:guaranteed-divisor}
\end{equation}
divides $J_F(z_1,z_2)$.
\end{proposition}

\begin{proof}
Each term of \eqref{eq:gap-ideal-expansion} has the valuation displayed inside the minimum.
The valuation of a sum is at least the minimum valuation of its summands.
\end{proof}

The functional in \eqref{eq:newton-floor} is the lower support function of the shifted Newton
polygon evaluated against a prime-dependent weight vector.  Unlike a determinant valuation,
it is a scalar certificate and can be optimized directly over the mask support.

\begin{criterion}[Product-formula mask criterion]
\label{crit:product-formula}
Suppose $F(z_1,z_2)\ne0$.  Then
\begin{equation}
 \mathcal D_F(z_1,z_2)
 \le |J_F(z_1,z_2)|
 \le
 \sum_{i,j}|\gamma_{ij}|
 |G_1|^i|G_2|^jS_1^{d-i}S_2^{e-j}.
 \label{eq:product-formula-sandwich}
\end{equation}
Therefore a contradiction follows if the explicit right-hand side is strictly smaller than
$\mathcal D_F(z_1,z_2)$.
\end{criterion}

\begin{proof}
The first inequality follows because the nonzero integer $J_F$ is divisible by
$\mathcal D_F$; the second is the triangle inequality applied to
\eqref{eq:gap-ideal-expansion}.
\end{proof}

\begin{example}[A positive cross-face mask]
\label{ex:positive-cross-face}
For $r\ge1$,
\begin{equation}
 Q_r(U,V)=(V+1)(U-1)^{2r}+(U+1)(V-1)^{2r}
 \label{eq:positive-cross-face}
\end{equation}
has contact order $2r$ and is strictly positive for positive $(U,V)\ne(1,1)$.  For $r=2$,
full bidegree-$(4,4)$ clearing gives the exact identity
\begin{align}
 S_1^4S_2^4Q_2(z_1,z_2)
 ={}&(R_2+S_2)S_2^3G_1^4
   +(R_1+S_1)S_1^3G_2^4.
 \label{eq:positive-cross-cleared}
\end{align}
This example guarantees nonvanishing without a determinant.  It will nevertheless fail at
private denominator primes unless the cross-face values $z_2+1$ and $z_1+1$ provide the
necessary divisibility; \cref{sec:leading-faces} makes that statement exact.
\end{example}

\begin{statusbox}{\statusproved{} for the gap-ideal theorem, the primewise Newton floor, and
the scalar product-formula criterion.  These statements are finite identities and valuation
inequalities, with no conjectural input.}
\end{statusbox}

\section{Leading faces: the private-prime prototype}
\label{sec:leading-faces}

The full clearing factor $S_1^dS_2^e$ is usually much larger than the reduced denominator of
$F(z_1,z_2)$, so a useful construction must force cancellation between that clearing factor
and the cleared numerator.  The next theorem identifies the first obstruction exactly.

Assume now that $z_1=R_1/S_1$ and $z_2=R_2/S_2$ are in lowest terms, and that $d\ge1$.  Write
\begin{equation}
 F(U,V)=\sum_{i=0}^d f_i(V)U^i,
 \qquad f_i\in\Z[V],
 \qquad \deg f_i\le e,
 \qquad f_d\ne0.
 \label{eq:U-slices}
\end{equation}
The polynomial $f_d(V)$ is the top $U$-face of the Newton polygon.  Put
\begin{equation}
 A_i=S_2^e f_i(z_2)\in\Z.
 \label{eq:Ai-cleared}
\end{equation}
Then
\begin{equation}
 J_F=S_1^dS_2^eF(z_1,z_2)
     =\sum_{i=0}^d A_iR_1^iS_1^{d-i}.
 \label{eq:leading-face-sum}
\end{equation}

\begin{theorem}[Leading-face survival at a private prime]
\label{thm:leading-face-survival}
Let $p\mid S_1$ and $p\nmid S_2$, and put $s=v_p(S_1)$.  If
\begin{equation}
 v_p(A_d)<s,
 \label{eq:leading-face-shallow}
\end{equation}
then
\begin{equation}
 v_p(J_F)=v_p(A_d).
 \label{eq:leading-face-exact-vp}
\end{equation}
Consequently the exponent of $p$ in the reduced denominator of $F(z_1,z_2)$ is exactly
\begin{equation}
 d s-v_p(A_d).
 \label{eq:reduced-den-private}
\end{equation}
In particular, if $p\nmid A_d$, the entire pole $p^{ds}$ survives.
\end{theorem}

\begin{proof}
Since $z_1$ is reduced, $p\nmid R_1$.  The $i=d$ term in
\eqref{eq:leading-face-sum} has valuation $v_p(A_d)$.  Every term with $i<d$ has valuation at
least $(d-i)s\ge s$, because $A_i$ is an integer.  Under
\eqref{eq:leading-face-shallow}, the top-face term is the unique term of minimal valuation,
so the ultrametric inequality is an equality and gives
\eqref{eq:leading-face-exact-vp}.  The full clearing denominator has $p$-valuation $ds$,
and $p\nmid S_2$, proving \eqref{eq:reduced-den-private}.
\end{proof}

\begin{corollary}[Necessary first-layer capture]
\label{cor:first-layer-capture}
To reduce the $p$-part of the denominator below $p^{(d-1)s}$ at a prime satisfying
$p\mid S_1$, $p\nmid S_2$, it is necessary that
\begin{equation}
 p^s\mid S_2^e f_d(z_2).
 \label{eq:first-layer-capture}
\end{equation}
If the divisibility holds at every prime in the private part
\begin{equation}
 P_1=\prod_{p\mid S_1,\ p\nmid S_2}p^{v_p(S_1)},
 \label{eq:private-part}
\end{equation}
then either $f_d(z_2)=0$ or
\begin{equation}
 \log P_1
 \le e h(z_2)+\log\|f_d\|_1,
 \label{eq:first-layer-height-budget}
\end{equation}
where $h(z_2)=\log\max(R_2,S_2)$ and $\|f_d\|_1$ is the coefficient $\ell^1$ norm.
\end{corollary}

\begin{proof}
The first statement is the contrapositive of \cref{thm:leading-face-survival}.  If it holds
for all private primes, then $P_1\mid A_d$.  Unless $A_d=0$,
\[
 P_1\le |A_d|=|S_2^e f_d(R_2/S_2)|
 \le \|f_d\|_1\max(R_2,S_2)^e,
\]
which gives \eqref{eq:first-layer-height-budget}.
\end{proof}

This is a genuinely asymmetric statement: cancellation of the top $U$-pole layer must be
paid for by the value of the top $U$-face at the other near-unit.  The symmetric theorem is
obtained by exchanging $U,V$.

\subsection{Monomial transport versus unit-root capture}
\label{subsec:monomial-vs-root}

Factor the top face as
\begin{equation}
 f_d(V)=V^m\widehat f_d(V),
 \qquad \widehat f_d(0)\ne0.
 \label{eq:face-factor}
\end{equation}
At a prime $p\nmid S_2$, the monomial $V^m$ contributes
$m v_p(R_2)$ to $v_p(A_d)$.  This is deterministic transport of valuation from the numerator
of $z_2$ to the denominator of $z_1$.  It is exactly the min-plus or tropical mechanism
already audited in the supplied report.

After that monomial contribution is removed, any additional valuation must come from the
congruence
\begin{equation}
 \widehat f_d(z_2)\equiv0\pmod{p^k}
 \label{eq:unit-root-capture}
\end{equation}
when $z_2$ is a $p$-adic unit.  We call this \emph{unit-root capture}.  It is the genuinely
new denominator mechanism isolated by the mask framework.

\begin{corollary}[Face-monic obstruction]
\label{cor:face-monic}
Suppose the top $U$-face is the constant $f_d(V)=\pm1$.  Then at every prime
$p\mid S_1$ with $p\nmid S_2$, the full $p^{d v_p(S_1)}$ denominator survives.  More
generally, if $f_d(V)=\pm V^m$ and $p\nmid R_2S_2$, the full pole survives.
\end{corollary}

\begin{proof}
In the first case $A_d=\pm S_2^e$; in the second it is a $p$-adic unit times a power of the
$p$-adic unit $z_2$.  Apply \cref{thm:leading-face-survival}.
\end{proof}

\begin{example}[Natural positive masks are face-monic]
The mask
\begin{equation}
 F_r(U,V)=(U-1)^{2r}+(V-1)^{2r}
 \label{eq:basic-positive-mask}
\end{equation}
is positive for positive $(U,V)\ne(1,1)$ and has contact order $2r$.  Its top $U$-face and
top $V$-face are both $1$.  Hence every denominator prime private to either near-unit
survives with full multiplicity.  Positivity solves nonvanishing but leaves the denominator
tax untouched.
\end{example}

\subsection{The pole-layer Newton polygon}

The first-layer theorem is the first edge of a more complete local picture.  At a prime
$p\mid S_1$, $p\nmid S_2$, put $s=v_p(S_1)$ and consider the points
\begin{equation}
 \bigl(i,v_p(A_i)\bigr),\qquad 0\le i\le d.
 \label{eq:pole-layer-points}
\end{equation}
The term with index $i$ in \eqref{eq:leading-face-sum} has valuation
$v_p(A_i)+(d-i)s$.  Thus $v_p(J_F)$ is at least the minimum of the affine functional
\begin{equation}
 \ell_{p,s}(i,t)=t+(d-i)s
 \label{eq:pole-functional}
\end{equation}
over the coefficient Newton polygon, with equality whenever the minimum is unique.  Full
cancellation of successive pole layers requires the right-hand vertices to rise by at least
$s$ per horizontal step, or requires exact cancellation among tied minimal terms.  Either
phenomenon translates to high $p$-adic divisibility of a short list of face evaluations.

\begin{target}[Iterated face-capture theorem]
\label{target:iterated-face}
Develop a global, coefficient-height-efficient bound for the total number of canceled pole
layers in terms of the evaluated face chain
\[
 f_d(z_2),f_{d-1}(z_2),\ldots
\]
and the resultants of successive faces.  The desired form is a lower bound for the reduced
denominator that remains valid even when several local minima tie and cancel.
\end{target}

The first-layer result is already enough to rule out all face-monic constructions and to
turn cyclotomic faces into an explicit finite-order covering problem.

\begin{statusbox}{\statusproved{} for leading-face survival, the private first-layer
condition, its height budget, and the face-monic obstruction.  \Cref{target:iterated-face} is
open and is one of the two strongest concrete follow-up problems identified here.}
\end{statusbox}

\section{Cyclotomic cross-faces: a finite private-prime prototype}
\label{sec:cyclotomic-faces}

A polynomial leading face can capture a unit-private prime only by having the other near-unit
as a root modulo that prime.  Cyclotomic polynomials give the most structured way to enforce
such roots while retaining positivity on the positive real axis.

For a finite set $\mathcal M\subset\{2,3,\ldots\}$ define
\begin{equation}
 A_{\mathcal M}(T)=\prod_{m\in\mathcal M}\Phi_m(T),
 \qquad
 E(\mathcal M)=\deg A_{\mathcal M}=\sum_{m\in\mathcal M}\varphi(m).
 \label{eq:cyclotomic-cover-poly}
\end{equation}

\begin{lemma}[Positivity and coefficient budget]
\label{lem:cyclotomic-positive}
For every $m\ge2$, $\Phi_m(t)>0$ for all $t>0$.  Moreover
\begin{equation}
 \|A_{\mathcal M}\|_1\le2^{E(\mathcal M)}.
 \label{eq:cyclotomic-l1}
\end{equation}
\end{lemma}

\begin{proof}
The polynomial $\Phi_m$ has no positive real root for $m\ge2$, has positive leading
coefficient, and has value $1$ at $0$, so it is positive on $(0,\infty)$.  Its roots all have
modulus $1$.  Therefore the absolute value of its coefficient of degree $j$ is at most the
corresponding binomial coefficient, and $\|\Phi_m\|_1\le2^{\varphi(m)}$.  The $\ell^1$ norm
is submultiplicative under polynomial multiplication.
\end{proof}

For $m\ge2$ define the homogenized cyclotomic polynomial by
\begin{equation}
 \Phi_m^{\mathrm{hom}}(R,S):=S^{\varphi(m)}\Phi_m(R/S)\in\Z[R,S].
 \label{eq:hom-cyclotomic}
\end{equation}

\begin{proposition}[Cyclotomic order detection]
\label{prop:cyclotomic-order}
Let $z=R/S$ be reduced and let $p\nmid RS\prod_{m\in\mathcal M}m$.  Then
\begin{equation}
 p\mid S^{E(\mathcal M)}A_{\mathcal M}(z)
 \quad\Longleftrightarrow\quad
 \ord_p(RS^{-1})\in\mathcal M,
 \label{eq:order-detection}
\end{equation}
where the order is in $\mathbb F_p^\times$.
\end{proposition}

\begin{proof}
The denominator is a $p$-adic unit.  For $p\nmid m$, an element of
$\mathbb F_p^\times$ is a root of $\Phi_m$ if and only if it has exact multiplicative order
$m$.  The product vanishes precisely when one factor does.
\end{proof}

For finite order sets $\mathcal M,\mathcal N$ and $r\ge1$, consider
\begin{equation}
 \mathcal F_{r,\mathcal M,\mathcal N}(U,V)
 =A_{\mathcal M}(V)(U-1)^{2r}
  +A_{\mathcal N}(U)(V-1)^{2r}.
 \label{eq:cyclotomic-face-mask}
\end{equation}

\begin{proposition}[Positive cyclotomic-face mask]
\label{prop:positive-cyclotomic-mask}
The mask in \eqref{eq:cyclotomic-face-mask} is strictly positive for all positive
$(U,V)\ne(1,1)$ and has contact order exactly $2r$.  If
$2r>E(\mathcal N)$, its top $U$-face is $A_{\mathcal M}(V)$; if
$2r>E(\mathcal M)$, its top $V$-face is $A_{\mathcal N}(U)$.
\end{proposition}

\begin{proof}
Both cyclotomic products are positive by \cref{lem:cyclotomic-positive}; the even powers are
nonnegative and cannot vanish simultaneously away from $(1,1)$.  Since
$A_{\mathcal M}(1),A_{\mathcal N}(1)>0$, the first nonzero homogeneous part has degree $2r$.
The face statements follow by comparing degrees.
\end{proof}

Combining this proposition with \cref{thm:leading-face-survival,prop:cyclotomic-order} gives
the following exact audit.

\begin{corollary}[Order-cover necessity]
\label{cor:order-cover-necessity}
Assume $2r>E(\mathcal N)$.  Let $p\mid S_1$, $p\nmid R_2S_2$, and
$p\nmid\prod_{m\in\mathcal M}m$.  For the first $p$-pole layer of the $U$-direction to
cancel in \eqref{eq:cyclotomic-face-mask}, it is necessary that
\begin{equation}
 \ord_p(z_2\bmod p)\in\mathcal M.
 \label{eq:order-cover-necessary}
\end{equation}
The symmetric condition holds for denominator primes of $z_2$.
\end{corollary}

\begin{proof}
The top $U$-face is $A_{\mathcal M}(V)$.  First-layer cancellation requires it to vanish
modulo $p$ by \cref{cor:first-layer-capture}; \cref{prop:cyclotomic-order} translates this to
\eqref{eq:order-cover-necessary}.
\end{proof}

\subsection{Depth, not merely residue order}

Covering a prime modulo $p$ cancels at most the first pole layer.  If
$v_p(S_1)=s$ is large, the full condition is
\begin{equation}
 v_p\bigl(S_2^{E(\mathcal M)}A_{\mathcal M}(z_2)\bigr)\ge s.
 \label{eq:deep-order-cover}
\end{equation}
For a single factor with $p\nmid m$, the left side is
$v_p(z_2^m-1)$ once the exact order is $m$.  Lifting-the-exponent arguments and effective
$p$-adic logarithm bounds then show that a fixed order factor usually supplies only
logarithmic depth in the exponent-vector height, while $s$ can grow linearly.  This is where
recent sharp bounds for two $p$-adic logarithms can be inserted.

\begin{target}[Finite-pair low-cost deep order cover]
\label{target:deep-order-cover}
For a family of finite or nonadjacent output near-unit pairs admitting a private denominator
prime, let $H$ be the maximum exponent-vector height.  Determine whether the private
denominator mass of the first coordinate can be covered to full $p$-adic depth by a product
$A_H$ evaluated at the second coordinate with
\begin{equation}
 \deg A_H=o(H),
 \qquad
 \log\normone{A_H}=o(H),
 \label{eq:low-cost-cover}
\end{equation}
and analogously in the reverse direction.  Monomial factors are to be removed
first, so the target concerns genuine unit-root capture.
\end{target}

The target is intentionally local and candidate-specific.  It can matter for finite or
nonadjacent configurations, but \cref{thm:support-stabilization} below shows that sufficiently
deep adjacent pairs have no private structural denominator primes.  The asymptotic endgame is
therefore the shared-pole problem developed in \cref{sec:stable-support}.

\subsection{Why the order-cover formulation is useful}

It separates four effects that are easily conflated.

\begin{enumerate}
\item \textbf{Toric aliases.}  These vanish identically and are removed before any contact is
  counted.
\item \textbf{Monomial valuation transport.}  This is a linear cone problem on the fixed
  valuation vectors of $M,N$ and belongs to the tropical regime already studied.
\item \textbf{Finite private-prime unit-root capture.}  This is the order-cover mechanism
  visible in the top faces of the cyclotomic prototype.
\item \textbf{Shared structural-face cancellation.}  For sufficiently deep adjacent pairs,
  the same structural primes divide both denominators; multi-term initial forms must then
  cancel repeatedly on the shared Newton fan.
\end{enumerate}

The decomposition is exact.  The third mechanism remains a useful finite diagnostic, while
the stable-support theorem below shows that the fourth is the asymptotic endgame.

\begin{statusbox}{\statusproved{} for positivity, coefficient budget, residue-order
detection, and the resulting necessary order-cover condition.  \Cref{target:deep-order-cover}
is an open finite-pair prototype.  It is superseded asymptotically by the shared structural-face
target of \cref{target:shared-face-capture}.}
\end{statusbox}

\section{Stable structural support and the shared-pole Newton fan}
\label{sec:stable-support}

The private-prime analysis above is a useful local model, but it is not the asymptotic
geometry of the actual adjacent continued-fraction near-units.  For those pairs the numerator
and denominator supports eventually stabilize.  Thus the two coordinates do not exchange
fresh denominator primes; they carry the same fixed structural primes at different linear
depths.  This changes the surviving endgame from private-prime order covering to simultaneous
cancellation on shared $p$-adic Newton faces.

Keep the convergent notation
\[
 \nu_k=(-p_k,q_k),\qquad
 z_k=\frac{N^{q_k}}{M^{p_k}}=\frac{R_k}{S_k}
 \quad\text{in lowest terms}.
\]
For a prime $\ell\mid MN$, put
\begin{equation}
 m_\ell=v_\ell(M),\qquad n_\ell=v_\ell(N),\qquad
 u_{\ell,k}=v_\ell(z_k)=q_kn_\ell-p_km_\ell.
 \label{eq:structural-valuation}
\end{equation}
Since $p_k/q_k\to\thetaAE$, one has the exact decomposition
\begin{equation}
 u_{\ell,k}
 =q_k(n_\ell-\thetaAE m_\ell)
  +m_\ell(q_k\thetaAE-p_k).
 \label{eq:valuation-slope-decomposition}
\end{equation}

\begin{theorem}[Stable structural support]
\label{thm:support-stabilization}
Define
\begin{equation}
 \mathcal P_+=\{\ell\mid MN:n_\ell-\thetaAE m_\ell>0\},
 \qquad
 \mathcal P_-=\{\ell\mid MN:n_\ell-\thetaAE m_\ell<0\}.
 \label{eq:stable-support-sets}
\end{equation}
Then $\mathcal P_+$ and $\mathcal P_-$ are disjoint and nonempty, and their union is
the set of primes dividing $MN$.  For all sufficiently large $k$,
\begin{equation}
 R_k=\prod_{\ell\in\mathcal P_+}\ell^{u_{\ell,k}},
 \qquad
 S_k=\prod_{\ell\in\mathcal P_-}\ell^{-u_{\ell,k}}.
 \label{eq:stable-reduced-support}
\end{equation}
Moreover, for every $\ell\mid MN$,
\begin{equation}
 |u_{\ell,k}|
 =q_k|n_\ell-\thetaAE m_\ell|+O_{M,N}(q_{k+1}^{-1}),
 \label{eq:linear-pole-depth}
\end{equation}
so each nonzero numerator or denominator depth is linear in $q_k$.
In particular, for all sufficiently large adjacent pairs $(z_k,z_{k+1})$ there is no prime
private to one reduced denominator: both denominators are supported on the same fixed set
$\mathcal P_-$.
\end{theorem}

\begin{proof}
If $\ell\mid MN$, then $(m_\ell,n_\ell)\ne(0,0)$.  The coefficient
$n_\ell-\thetaAE m_\ell$ cannot vanish, because otherwise either $m_\ell=0=n_\ell$ or
$\thetaAE=n_\ell/m_\ell\in\Q$, contradicting the irrationality of $\log3/\log2$.
The convergent estimate $|q_k\thetaAE-p_k|<q_{k+1}^{-1}$ in
\eqref{eq:valuation-slope-decomposition} therefore fixes the sign for all sufficiently large
$k$ and gives \eqref{eq:linear-pole-depth}.  Reduced numerator and denominator valuations are
the positive and negative parts of $u_{\ell,k}$, proving
\eqref{eq:stable-reduced-support}.

Finally,
\[
 \sum_{\ell\mid MN}(n_\ell-\thetaAE m_\ell)\log\ell
 =\log N-\thetaAE\log M=0.
\]
Every summand coefficient is nonzero.  Hence they cannot all have the same sign, so both
$\mathcal P_+$ and $\mathcal P_-$ are nonempty.
\end{proof}

The theorem is easy to miss because $z_k$ itself tends to $1$ and alternates around it.  The
near-cancellation is not caused by small prime valuations.  It is a cancellation between two
fixed, linearly growing valuation profiles whose logarithmic weights almost balance.

\subsection{Adjacent valuation transport}

Let
\begin{equation}
 \varepsilon_k=p_{k+1}q_k-p_kq_{k+1}\in\{\pm1\}.
 \label{eq:convergent-orientation}
\end{equation}
The unimodularity of adjacent convergents gives exact transport identities at every
structural prime.

\begin{proposition}[Adjacent valuation transport and toric balance]
\label{prop:adjacent-valuation-transport}
For every prime $\ell\mid MN$,
\begin{align}
 q_{k+1}u_{\ell,k}-q_ku_{\ell,k+1}&=\varepsilon_k m_\ell,
 \label{eq:q-valuation-transport}\\
 p_{k+1}u_{\ell,k}-p_ku_{\ell,k+1}&=\varepsilon_k n_\ell.
 \label{eq:p-valuation-transport}
\end{align}
Equivalently, in $\Q_{>0}^{\times}$,
\begin{align}
 \frac{z_k^{q_{k+1}}}{z_{k+1}^{q_k}}&=M^{\varepsilon_k},
 \label{eq:q-toric-balance}\\
 \frac{z_k^{p_{k+1}}}{z_{k+1}^{p_k}}&=N^{\varepsilon_k}.
 \label{eq:p-toric-balance}
\end{align}
\end{proposition}

\begin{proof}
Insert $u_{\ell,j}=q_jn_\ell-p_jm_\ell$ and use
$p_{k+1}q_k-p_kq_{k+1}=\varepsilon_k$.  The multiplicative identities follow by the same
exponent calculation before taking valuations.
\end{proof}

Thus the two shared pole depths are not independent random data.  Their large linear parts
are tied by a determinant-one defect equal to one fixed valuation of $M$ or $N$.  This exact
coupling should be incorporated into any search rather than treated as an accidental
correlation.

\subsection{The shared-pole Newton polygon}

Fix a sufficiently large adjacent pair and abbreviate
\[
 z_1=z_k=R_1/S_1,\qquad z_2=z_{k+1}=R_2/S_2.
\]
Let
\begin{equation}
 F(U,V)=\sum_{\substack{0\le i\le d\\0\le j\le e}}c_{ij}U^iV^j\in\Z[U,V],
 \qquad
 J_F=S_1^dS_2^eF(z_1,z_2)\in\Z.
 \label{eq:shared-cleared-mask}
\end{equation}
For $\ell\in\mathcal P_-$ write $s_{\ell,j}=-u_{\ell,j}=v_\ell(S_j)$ for every
sufficiently large $j$.  For the fixed adjacent pair abbreviate
\begin{equation}
 s_{\ell,1}=s_{\ell,k}=v_\ell(S_1),
 \qquad
 s_{\ell,2}=s_{\ell,k+1}=v_\ell(S_2).
 \label{eq:shared-pole-depths}
\end{equation}
Both depths are positive and of order $q_k$.

\begin{theorem}[Exact shared-pole Newton law]
\label{thm:shared-pole-newton}
For $\ell\in\mathcal P_-$ define
\begin{equation}
 \lambda_{ij}^{(\ell)}
 =v_\ell(c_{ij})+(d-i)s_{\ell,1}+(e-j)s_{\ell,2}
 \label{eq:shared-newton-weight}
\end{equation}
for $c_{ij}\ne0$, and let $\lambda_\ell=\min\lambda_{ij}^{(\ell)}$.  Then
\begin{equation}
 v_\ell(J_F)\ge\lambda_\ell.
 \label{eq:shared-newton-lower}
\end{equation}
If the minimum is attained at a unique exponent pair, equality holds.
More generally, write $S_t=\ell^{s_{\ell,t}}S_t'$ with $\ell\nmid S_t'$.  After division by
$\ell^{\lambda_\ell}$, the residue of $J_F$ modulo $\ell$ is the sum of the residues of the
terms on the minimizing face:
\begin{equation}
 \operatorname{in}_{\ell}(F;z_1,z_2)
 =\sum_{\lambda_{ij}^{(\ell)}=\lambda_\ell}
 \overline{\ell^{-v_\ell(c_{ij})}c_{ij}}
 \overline{R_1}^{\,i}\overline{R_2}^{\,j}
 \overline{S_1'}^{\,d-i}\overline{S_2'}^{\,e-j}
 \in\mathbb F_\ell.
 \label{eq:shared-initial-form}
\end{equation}
Consequently
\begin{equation}
 v_\ell(J_F)>\lambda_\ell
 \quad\Longleftrightarrow\quad
 \operatorname{in}_{\ell}(F;z_1,z_2)=0.
 \label{eq:initial-form-cancellation}
\end{equation}
Iterating after division by the current power of $\ell$ gives an exact finite-depth
cancellation cascade.
\end{theorem}

\begin{proof}
Expand
\[
 J_F=\sum_{i,j}c_{ij}R_1^iR_2^jS_1^{d-i}S_2^{e-j}.
\]
Because $z_1,z_2$ are reduced and $\ell$ divides both denominators,
$\ell\nmid R_1R_2S_1'S_2'$.  The valuation of each summand is exactly
\eqref{eq:shared-newton-weight}.  The ultrametric inequality proves
\eqref{eq:shared-newton-lower}; a unique minimum cannot cancel.  Factoring the common power
$\ell^{\lambda_\ell}$ gives \eqref{eq:shared-initial-form} and
\eqref{eq:initial-form-cancellation}.
\end{proof}

\begin{corollary}[Unit-corner survival]
\label{cor:unit-corner-survival}
If $c_{de}$ is an $\ell$-adic unit, then $(d,e)$ is the unique minimizing exponent pair,
$v_\ell(J_F)=0$, and the full pole
\begin{equation}
 \ell^{d s_{\ell,1}+e s_{\ell,2}}
 \label{eq:full-shared-pole}
\end{equation}
survives in the reduced denominator of $F(z_1,z_2)$.
\end{corollary}

This is the shared-prime analogue of the earlier face-monic obstruction.  It also shows why
removing the top-right monomial is only the beginning: once the minimizing set has more than
one term, one must force its initial form to vanish at every stable denominator prime, then
repeat the argument at every higher $\ell$-adic layer.

\begin{example}[The positive two-axis mask]
\label{ex:positive-shared-mask}
For
\[
 F_r(U,V)=(U-1)^{2r}+(V-1)^{2r}
\]
one has
\begin{equation}
 J_{F_r}=G_1^{2r}S_2^{2r}+G_2^{2r}S_1^{2r}.
 \label{eq:positive-shared-clearing}
\end{equation}
At $\ell\in\mathcal P_-$ the two displayed terms have valuations
$2rs_{\ell,2}$ and $2rs_{\ell,1}$.  If the depths are unequal, the smaller one is unique and
\begin{equation}
 v_\ell(J_{F_r})=2r\min(s_{\ell,1},s_{\ell,2}),
 \label{eq:positive-shared-exact}
\end{equation}
so the reduced denominator retains
$\ell^{2r\max(s_{\ell,1},s_{\ell,2})}$.  If the depths tie, cancellation is possible only
through the explicit two-term residue in \eqref{eq:positive-shared-clearing}.  Positivity at
the real place therefore does not by itself buy denominator compression.
\end{example}

\begin{lemma}[Two-term contact cannot vanish at the output]
\label{lem:balanced-binomial-collapse}
Let $z_1,z_2$ be the adjacent output near-units.  A nonzero Laurent polynomial with at most
two monomials cannot satisfy both $F(1,1)=0$ and $F(z_1,z_2)=0$.
\end{lemma}

\begin{proof}
The condition at $(1,1)$ makes the two coefficients opposite, so after multiplying by a
monomial one has $F=c(U^aV^b-1)$ with $(a,b)\ne(0,0)$.  Vanishing at $(z_1,z_2)$ would give
$z_1^az_2^b=1$.  Adjacent relation vectors form a unimodular basis of $\Z^2$, so
$z_1,z_2$ generate the same rank-two group as $M,N$ and are multiplicatively independent.
Thus $a=b=0$, a contradiction.
\end{proof}

The lemma does not forbid a two-term initial form from canceling modulo one prime.  It says
that no single binomial supplies exact cancellation at every depth and every place.  A viable
mask must use genuinely multi-term, prime-dependent Newton faces.

\begin{target}[Shared structural-face capture]
\label{target:shared-face-capture}
Construct, for some sufficiently large adjacent pair, a nonzero mask $F_k$ with contact
$r_k$ such that:
\begin{enumerate}
\item its shifted coefficient norm is subheight and its contact supplies the archimedean gain
  required by the scalar criterion of \cref{sec:endgame};
\item for every $\ell\in\mathcal P_-$, the successive initial forms in
  \eqref{eq:shared-initial-form} vanish deeply enough that
  \begin{equation}
   \sum_{\ell\in\mathcal P_-}
   \bigl(d_ks_{\ell,k}+e_ks_{\ell,k+1}-v_\ell(J_{F_k})\bigr)_+\log\ell
   =O_x((d_k+e_k)\log q_k);
   \label{eq:full-shared-face-capture}
  \end{equation}
\item the mask is certified nonzero at $(z_k,z_{k+1})$ and is non-toric.
\end{enumerate}
For a fixed hypothetical candidate, one mask satisfying these finite conditions and the
strict scalar inequality contradicts that candidate.  A construction that works for every
hypothetical candidate would prove the \AEname\ conjecture.
\end{target}

\begin{statusbox}{\statusproved{} for stable support, adjacent valuation transport, the
shared-pole Newton law, unit-corner survival, the positive-mask calibration, and the two-term
nonvanishing lemma.  \Cref{target:shared-face-capture} is the new open endgame.  The earlier
private-prime and cyclotomic sections remain valid finite diagnostics, but
\cref{thm:support-stabilization} shows that they are not the asymptotic adjacent-convergent
mechanism.}
\end{statusbox}

\section{Three quantitative barriers for naive mask constructions}
\label{sec:quantitative-barriers}

The preceding sections identify the exact arithmetic objects.  We now combine them with the
continued-fraction scale.  The conclusions are negative, but they are useful negative results:
they rule out three broad architectures while leaving one sharply defined escape route.

Let $p_n/q_n$ be the continued-fraction convergents to $\thetaAE$, and put
\begin{equation}
 \nu_n=(-p_n,q_n),\qquad
 \delta_n=q_n\log3-p_n\log2,\qquad
 z_n=\sigma_{\nu_n}=\frac{N^{q_n}}{M^{p_n}}=e^{x\delta_n}.
 \label{eq:convergent-near-units}
\end{equation}
Write $z_n=R_n/S_n$ in lowest terms and set
\begin{equation}
 \Delta_n=\max\{|x\delta_n|,|x\delta_{n+1}|\}.
 \label{eq:adjacent-output-error}
\end{equation}
For all sufficiently large $n$, $\Delta_n<1/4$.

We use one standard external input.  An effective lower bound for linear forms in
$\log2,\log3$ gives constants $C>0$ and $\kappa>1$ such that
\begin{equation}
 |q\thetaAE-p|\ge Cq^{-\kappa}
 \qquad(p\in\Z,\ q\ge1).
 \label{eq:effective-two-log}
\end{equation}
For convergents this implies $q_{n+1}\ll q_n^{\kappa}$ and hence
\begin{equation}
 \log\frac1{\Delta_n}=O_x(\log q_n).
 \label{eq:delta-logarithmic}
\end{equation}
Only the existence of some effective exponent is needed; no numerical optimization enters.

\subsection{The full-clearing Taylor envelope cannot become subunit}

Let $F_n\in\Z[U,V]$ be nonzero, of bidegree at most $(d_n,e_n)$ and contact order
$r_n$ at $(1,1)$.  Write
\[
 \widetilde F_n(X,Y)=F_n(1+X,1+Y).
\]
Clearing the visible denominators and applying \cref{lem:contact-estimate} gives
\begin{equation}
 |J_n|
 \le S_n^{d_n}S_{n+1}^{e_n}
       \normone{\widetilde F_n}(2\Delta_n)^{r_n},
 \qquad
 J_n:=S_n^{d_n}S_{n+1}^{e_n}F_n(z_n,z_{n+1})\in\Z.
 \label{eq:full-clearing-envelope}
\end{equation}

\begin{theorem}[Full-clearing envelope barrier]
\label{thm:full-clearing-barrier}
Assume $d_n+e_n>0$.  Then the right side of
\eqref{eq:full-clearing-envelope} is greater than $1$ for all sufficiently large $n$,
uniformly over all choices of $F_n$.  Consequently the ordinary argument
\[
 \text{clear every visible denominator}+\text{bound by Taylor contact}
\]
can never prove $J_n=0$ or produce a nonzero integer of absolute value below one.
\end{theorem}

\begin{proof}
By \cref{prop:height-law}, norm equivalence, and $p_n\asymp q_n$, there is a constant
$c_x>0$ such that
\[
 \log S_n\ge c_xq_n-O_x(1),
 \qquad
 \log S_{n+1}\ge c_xq_{n+1}-O_x(1)
\]
for all large $n$.  The contact ceiling \cref{thm:contact-ceilings} gives
$r_n\le d_n+e_n$, while $\normone{\widetilde F_n}\ge1$.  Therefore the logarithm of the
right side in \eqref{eq:full-clearing-envelope} is at least
\[
 c_xq_n(d_n+e_n)-O_x(d_n+e_n)
 -(d_n+e_n)O_x(\log q_n),
\]
where \eqref{eq:delta-logarithmic} was used in the last term.  This is positive for all
sufficiently large $n$.
\end{proof}

\begin{remark}[What the theorem does not say]
The theorem does not lower-bound the actual value $|J_n|$ by the displayed envelope.  It says
that this particular triangle-inequality proof cannot make the integer subunit.  The two
possible escapes are cancellation inside the archimedean sum, or cancellation in the
\emph{reduced} denominator before clearing.  The first is an auxiliary-function problem; the second is the shared-pole
Newton-face problem isolated in \cref{sec:stable-support}.
\end{remark}

\subsection{Imported gap transport and the new shared-gap consequence}

For a relation vector $\nu=(-p,q)$ put
\[
 z_\nu=\frac{N^q}{M^p}=\frac{R_\nu}{S_\nu},
 \qquad
 G_\nu=R_\nu-S_\nu,
\]
where the fraction is reduced.  The supplied report already proves the determinant-transport
law for pairs of binomial gaps.  Before importing it, there is a useful reduced-gap
observation that removes the usual exceptional support.

\begin{lemma}[Structural primes cannot be shared by independent reduced gaps]
\label{lem:structural-primes-not-shared}
Let $\nu,\omega\in\Z^2$ be linearly independent.  Then
\begin{equation}
 \gcd(G_\nu,G_\omega)\ \text{is coprime to }MN.
 \label{eq:structural-primes-not-shared}
\end{equation}
\end{lemma}

\begin{proof}
Suppose that a prime $\ell\mid MN$ divides $G_\nu$.  Since
$\gcd(R_\nu,S_\nu)=1$, a prime dividing $R_\nu-S_\nu$ divides neither
$R_\nu$ nor $S_\nu$.  Hence
\[
 v_\ell(z_\nu)=0,
 \qquad
 \nu\cdot\bigl(v_\ell(M),v_\ell(N)\bigr)=0.
\]
If the same prime divides $G_\omega$, then $\omega$ is orthogonal to the same nonzero
valuation vector.  Its kernel in $\Q^2$ is one-dimensional, forcing $\nu$ and $\omega$ to be
linearly dependent, a contradiction.
\end{proof}

\begin{lemma}[Inherited determinant transport for reduced gaps]
\label{lem:reduced-gap-transport}
Let $p,q,r,s\in\N$ and let $\nu=(-p,q)$ and $\omega=(-r,s)$ be
independent, with $D=|ps-rq|$.  Then
\begin{equation}
 \gcd(G_\nu,G_\omega)
 \mid \gcd(M^D-1,N^D-1).
 \label{eq:reduced-gap-transport}
\end{equation}
\end{lemma}

\begin{proof}
This is the gap-determinant theorem proved in the supplied report, specialized to reduced
output gaps.  For completeness of the specialization, a prime power dividing both reduced
gaps is coprime to $MN$ by \cref{lem:structural-primes-not-shared}; it therefore divides the
two unreduced gaps $N^q-M^p$ and $N^s-M^r$.  The inherited determinant transport gives
$M^D\equiv N^D\equiv1$ modulo that prime power, proving
\eqref{eq:reduced-gap-transport} prime by prime.
\end{proof}

For adjacent convergents, $D=1$.  Thus
\begin{equation}
 \gcd(G_n,G_{n+1})\mid\gcd(M-1,N-1),
 \label{eq:adjacent-reduced-gap-fixed}
\end{equation}
with no exceptional structural factor.  Combined with the new gap-ideal divisor
\eqref{eq:gcd-gap-floor}, this contributes at most $\exp(O_x(r_n))$, whereas the visible
denominator scale is $\exp(\Omega_x(q_n(d_n+e_n)))$.  Since
$r_n\le d_n+e_n$, adjacent shared-gap amplification is exponentially too small.

The same mask-level conclusion survives nonadjacent pairs.

\begin{corollary}[Shared-gap mask barrier]
\label{cor:nonadjacent-gcd-barrier}
Consider independent pairs $\nu_j,\omega_j\in\Z^2$ with
\[
 H_j=\normone{\nu_j}+\normone{\omega_j}\longrightarrow\infty,
 \qquad
 \max(|\delta_{\nu_j}|,|\delta_{\omega_j}|)\longrightarrow0.
\]
Let $g_j=\gcd(G_{\nu_j},G_{\omega_j})$ and let $r_j\ge1$ be any mask contact order.  Then
\begin{equation}
 \log g_j=o(H_j),
 \qquad
 \log(g_j^{r_j})=o(H_jr_j).
 \label{eq:fresh-gcd-subheight}
\end{equation}
Consequently the universal shared-gap divisor supplied by contact cannot pay a denominator
tax of order $H_j(d_j+e_j)$ when $r_j\le d_j+e_j$.
\end{corollary}

\begin{proof}
Put $D_j=|\det(\nu_j,\omega_j)|$.  The exact identities
\eqref{eq:uncertainty-identities} give
$D_j\ll H_j\max(|\delta_{\nu_j}|,|\delta_{\omega_j}|)=o(H_j)$.
By \cref{lem:reduced-gap-transport},
$g_j\mid\gcd(M^{D_j}-1,N^{D_j}-1)$.  If $D_j$ is bounded, then $g_j$ is bounded.  If
$D_j\to\infty$, the theorem of Bugeaud--Corvaja--Zannier gives
\[
 \log\gcd(M^{D_j}-1,N^{D_j}-1)=o(D_j),
\]
because $M,N$ are multiplicatively independent.  Thus $\log g_j=o(H_j)$ in either case;
multiplication by $r_j$ proves the second estimate.
\end{proof}

\subsection{The private-prime obstruction is local; the shared-pole obstruction is asymptotic}

For finite or nonadjacent pairs that possess a private denominator prime, the face-monic
obstruction of \cref{cor:face-monic} remains exact.  It eliminates powers of one-variable
gaps, masks with unit or monomial top faces, and generic sparse masks lacking explicit
opposite-coordinate order control.

For sufficiently deep adjacent convergents, however, \cref{thm:support-stabilization} shows
that there are no private structural denominator primes at all.  The obstruction changes
form: at every $\ell\in\mathcal P_-$ both coordinates have poles of linear depth, and the
minimizing terms are selected by the shared weight
\[
 (i,j)\longmapsto v_\ell(c_{ij})+(d-i)s_{\ell,k}+(e-j)s_{\ell,k+1}.
\]
A unit top-right corner leaves the whole pole, while a missing corner creates a slanted
Newton face whose initial form must vanish.  Near-total compression requires this vanishing
to persist through almost the complete depth at every prime in the fixed set
$\mathcal P_-$.  Thus the three quantitative barriers leave one architecture: multi-term,
prime-specific cancellation on shared Newton faces, coupled by the exact transport
identities \eqref{eq:q-valuation-transport}--\eqref{eq:p-valuation-transport}.

\begin{statusbox}{\statusproved{} for the full-clearing envelope barrier, structural-prime lemma, and
shared-gap mask corollary; the determinant-transport lemma is inherited from the supplied report; \statusexternal{} for the effective two-logarithm estimate and the
Bugeaud--Corvaja--Zannier gcd bound.  The conclusions concern the specified proof
architectures, not the truth of the \AEname\ conjecture.}
\end{statusbox}

\section{The reduced-denominator endgame}
\label{sec:endgame}

The preceding no-go theorem used the visible denominator
$S_n^{d_n}S_{n+1}^{e_n}$.  That is deliberately crude.  The actual arithmetic quantity is
the reduced denominator of the rational mask value.

Let $F\in\Z[U,V]$ have contact order $r$, and suppose
\[
 F(z_1,z_2)=\frac{a_F}{b_F}
 \quad\text{in lowest terms},
 \qquad b_F>0.
\]
Set $\eta_i=\log z_i$ and $\Delta=\max(|\eta_1|,|\eta_2|)$.

\begin{criterion}[Exact scalar mask criterion]
\label{crit:scalar-mask}
Assume $0<\Delta\le1/4$ and $F(z_1,z_2)\ne0$.  If
\begin{equation}
 r\log\frac1{2\Delta}
 >\log b_F+\log\normone{F(1+X,1+Y)},
 \label{eq:scalar-mask-criterion}
\end{equation}
then no hypothetical counterexample producing $z_1,z_2$ can exist.
\end{criterion}

\begin{proof}
A nonzero reduced rational number satisfies $|F(z_1,z_2)|\ge1/b_F$.  On the other hand,
\cref{lem:contact-estimate} gives
\[
 |F(z_1,z_2)|
 \le\normone{F(1+X,1+Y)}(2\Delta)^r.
\]
Condition \eqref{eq:scalar-mask-criterion} makes the latter strictly smaller than $1/b_F$.
\end{proof}

This elementary criterion is the precise endgame.  It neither assumes nor hides a determinant.
Every candidate construction can be audited by three numbers: contact $r$, shifted coefficient
norm, and the \emph{actual reduced denominator} $b_F$.

\subsection{The denominator-compression index}

If $F$ has bidegree at most $(d,e)$, then $b_F$ divides the visible denominator
$B_F=S_1^dS_2^e$.  Define the integer compression factor and normalized compression index
\begin{equation}
 K_F=\frac{B_F}{b_F},
 \qquad
 \mathfrak c(F;z_1,z_2)=
 \frac{\log K_F}{d\log S_1+e\log S_2},
 \label{eq:compression-index}
\end{equation}
with the convention $\mathfrak c=0$ when $B_F=1$.  Thus $0\le\mathfrak c\le1$.
Criterion \eqref{eq:scalar-mask-criterion} becomes
\begin{equation}
 \log K_F>
 d\log S_1+e\log S_2
 +\log\normone{F(1+X,1+Y)}
 -r\log\frac1{2\Delta}.
 \label{eq:compression-threshold}
\end{equation}

\begin{theorem}[Near-total compression is necessary on adjacent convergents]
\label{thm:near-total-compression}
Let $F_n$ be masks on $(z_n,z_{n+1})$ with bidegrees $(d_n,e_n)$, contact $r_n$, and
$F_n(z_n,z_{n+1})\ne0$.  Assume
\begin{equation}
 \log\normone{F_n(1+X,1+Y)}=o\bigl(q_n(d_n+e_n)\bigr).
 \label{eq:subheight-coefficients}
\end{equation}
If the scalar criterion succeeds for infinitely many $n$, then
\begin{equation}
 \mathfrak c(F_n;z_n,z_{n+1})=1-o(1)
 \label{eq:compression-to-one}
\end{equation}
along that subsequence.  More quantitatively, the uncanceled denominator obeys
\begin{equation}
 \log b_{F_n}
 \le O_x\bigl((d_n+e_n)\log q_n\bigr)
   -\log\normone{F_n(1+X,1+Y)}.
 \label{eq:uncanceled-polyheight}
\end{equation}
\end{theorem}

\begin{proof}
The contact ceiling gives $r_n\le d_n+e_n$, and
\eqref{eq:delta-logarithmic} gives
$r_n\log(1/(2\Delta_n))=O_x((d_n+e_n)\log q_n)$.  Criterion
\eqref{eq:scalar-mask-criterion} then implies \eqref{eq:uncanceled-polyheight}.  By
\cref{prop:height-law}, the visible denominator has logarithm
$\Omega_x(q_n(d_n+e_n))$.  Under \eqref{eq:subheight-coefficients}, subtracting
\eqref{eq:uncanceled-polyheight} from $\log B_{F_n}$ gives
$\log K_{F_n}=(1-o(1))\log B_{F_n}$.
\end{proof}

Thus ``some cancellation'' is not enough.  The mask must remove all but a polynomial-height
fraction of an exponentially large visible denominator.  Stable structural support and the
shared-pole Newton law make that requirement exact prime by prime.

\subsection{Exact primewise reduced-denominator formula}

For all sufficiently large $n$, every prime divisor of $S_nS_{n+1}$ lies in the fixed set
$\mathcal P_-$.  Put
\[
 J_{F_n}=S_n^{d_n}S_{n+1}^{e_n}F_n(z_n,z_{n+1})
\]
and, for $\ell\in\mathcal P_-$, define the uncanceled pole depth
\begin{equation}
 \rho_{\ell,n}(F_n)
 =\bigl(d_ns_{\ell,n}+e_ns_{\ell,n+1}-v_\ell(J_{F_n})\bigr)_+.
 \label{eq:uncanceled-shared-depth}
\end{equation}

\begin{proposition}[Exact shared-prime denominator law]
\label{prop:exact-shared-denominator}
If $F_n(z_n,z_{n+1})\ne0$ and $b_{F_n}$ is its reduced denominator, then
\begin{equation}
 b_{F_n}=\prod_{\ell\in\mathcal P_-}\ell^{\rho_{\ell,n}(F_n)},
 \qquad
 \log b_{F_n}
 =\sum_{\ell\in\mathcal P_-}\rho_{\ell,n}(F_n)\log\ell.
 \label{eq:exact-shared-denominator}
\end{equation}
\end{proposition}

\begin{proof}
The unreduced rational value is $J_{F_n}/(S_n^{d_n}S_{n+1}^{e_n})$.  At a prime
$\ell\in\mathcal P_-$ the denominator has valuation
$d_ns_{\ell,n}+e_ns_{\ell,n+1}$; reduction against the numerator $J_{F_n}$ leaves exactly
\eqref{eq:uncanceled-shared-depth}.  No other prime occurs in the visible denominator.
\end{proof}

\begin{corollary}[Scalar success forces simultaneous shared-face capture]
\label{cor:primewise-capture-forced}
Assume the subheight coefficient condition \eqref{eq:subheight-coefficients}.  If the scalar
criterion succeeds along infinitely many adjacent pairs, then
\begin{equation}
 \sum_{\ell\in\mathcal P_-}\rho_{\ell,n}(F_n)\log\ell
 =O_x((d_n+e_n)\log q_n)
 \label{eq:shared-residual-polyheight}
\end{equation}
along that subsequence.  Since $\mathcal P_-$ is fixed and every visible pole depth is
$\Theta_x(q_n)$, the mask must capture a $1-o(1)$ fraction of the pole depth at every stable
denominator prime that contributes a positive proportion of the visible height.
\end{corollary}

\begin{proof}
Combine \cref{prop:exact-shared-denominator} with
\eqref{eq:uncanceled-polyheight}.  The final assertion follows from
\eqref{eq:linear-pole-depth} and the finiteness of $\mathcal P_-$.  A primewise residual may
be redistributed among the finitely many primes, but the weighted total is only
polylogarithmic while each uncanceled linear fraction would contribute $\Theta(q_n)$.
\end{proof}

This is the precise surviving endgame.  At each stable denominator prime, the Newton initial
form \eqref{eq:shared-initial-form} must cancel; after division by the first canceled layer,
the next initial form must cancel; and so on through almost the complete linear depth.  The
exact transport laws show that the two depth vectors differ by a determinant-one defect, so
the required cancellations are strongly coupled across adjacent indices.

\subsection{Constructive and obstructive targets}

The constructive target is \cref{target:shared-face-capture}.  Cyclotomic faces remain a
useful finite source of residue identities, but stable support means that a successful mask
must combine them into multi-term shared faces rather than assign a private prime to the
opposite coordinate.  Positivity, a Sturm certificate, or another exact argument must also
keep the final rational value nonzero.

The complementary obstruction target is to prove that every nonzero, non-toric mask of
contact $r$ and subheight shifted coefficient norm leaves a positive proportion of the
shared pole mass uncanceled:
\begin{equation}
 \log b_F\ge c_x(d\log S_1+e\log S_2)-O_x(r\log H)
 \label{eq:shared-face-lower-target}
\end{equation}
for some $c_x>0$, where $H$ is the exponent-vector height.  Even a fixed positive $c_x$
would contradict the necessary polyheight bound \eqref{eq:uncanceled-polyheight}.  By
\cref{thm:shared-pole-newton}, this is a finite-prime theorem about repeated cancellation on
shared Newton faces, not an archimedean approximation theorem.

\begin{statusbox}{\statusproved{} for the scalar criterion, compression reformulation,
near-total-compression theorem, exact shared-prime denominator law, and forced simultaneous
capture.  \Cref{target:shared-face-capture} and \eqref{eq:shared-face-lower-target} are the
constructive and obstructive open endgames.}
\end{statusbox}

\section{Exact and numerical reconnaissance}
\label{sec:reconnaissance}

The computations in this section are diagnostic.  None is used as an unproved premise in the
theorems above.  Their purpose is to test constants, catch false high-contact masks that are
actually toric, and identify the scale a constructive search must beat.

\subsection{Continued-fraction near-unit bases}

The first convergents to $\thetaAE=\log3/\log2$ and their logarithmic errors are shown below.
The last column is
\begin{equation}
 \mathcal U_n=
 \frac{2H_n\max(|\delta_{n-1}|,|\delta_n|)}{\log3},
 \qquad
 H_n=\max(p_{n-1},q_{n-1},p_n,q_n).
 \label{eq:uncertainty-ratio}
\end{equation}
The determinant-one identity forces $\mathcal U_n\ge1$.  Values close to one indicate that the
simple uncertainty bound is already near the observed scale.

\begin{center}
\small
\begin{tabular}{r r r r r}
\toprule
$n$ & $p_n/q_n$ & $\delta_n=q_n\log3-p_n\log2$ & $q_n|\delta_n|$ & $\mathcal U_n$ \\
\midrule
0  & $1/1$          & $ 4.0546510811\!\times10^{-1}$ & $4.0547\!\times10^{-1}$ & -- \\
1  & $2/1$          & $-2.8768207245\!\times10^{-1}$ & $2.8768\!\times10^{-1}$ & 1.4763 \\
2  & $3/2$          & $ 1.1778303566\!\times10^{-1}$ & $2.3557\!\times10^{-1}$ & 1.5711 \\
3  & $8/5$          & $-5.2116001139\!\times10^{-2}$ & $2.6058\!\times10^{-1}$ & 1.7155 \\
4  & $19/12$        & $ 1.3551033379\!\times10^{-2}$ & $1.6261\!\times10^{-1}$ & 1.8016 \\
5  & $65/41$        & $-1.1462901003\!\times10^{-2}$ & $4.6998\!\times10^{-1}$ & 1.6037 \\
6  & $84/53$        & $ 2.0881323758\!\times10^{-3}$ & $1.1067\!\times10^{-1}$ & 1.7532 \\
7  & $485/306$      & $-1.0222391238\!\times10^{-3}$ & $3.1281\!\times10^{-1}$ & 1.8437 \\
8  & $1054/665$     & $ 4.3654128136\!\times10^{-5}$ & $2.9030\!\times10^{-2}$ & 1.9615 \\
9  & $24727/15601$  & $-1.8194176704\!\times10^{-5}$ & $2.8385\!\times10^{-1}$ & 1.9657 \\
10 & $50508/31867$  & $ 7.2657747287\!\times10^{-6}$ & $2.3153\!\times10^{-1}$ & 1.6355 \\
11 & $125743/79335$ & $-3.6626272464\!\times10^{-6}$ & $2.9059\!\times10^{-1}$ & 1.6630 \\
12 & $176251/111202$& $ 3.6031474823\!\times10^{-6}$ & $4.0068\!\times10^{-1}$ & 1.1559 \\
13 & $301994/190537$& $-5.9489764120\!\times10^{-8}$ & $1.1334\!\times10^{-2}$ & 1.9807 \\
\bottomrule
\end{tabular}
\end{center}

The large partial quotient $55$ following the thirteenth convergent creates a spectacularly
small single error, but the adjacent pair still obeys the determinant-one floor.  This is the
concrete form of the geometry-of-numbers obstruction: a single exceptional direction does not
create two independent exceptional directions.

\subsection{Symbolic mask checks}

The univariate Thue--Morse masks
\[
 P_r(T)=\prod_{j=0}^{r-1}(1-T^{2^j})
\]
have the following exact data.

\begin{center}
\small
\begin{tabular}{r r r r r}
\toprule
$r$ & degree & support & coefficient $\ell^1$ norm & multiplicity at $1$ \\
\midrule
1&1&2&2&1\\
2&3&4&4&2\\
3&7&8&8&3\\
4&15&16&16&4\\
5&31&32&32&5\\
6&63&64&64&6\\
7&127&128&128&7\\
\bottomrule
\end{tabular}
\end{center}

The growth is attractive from a moment-cancellation viewpoint, but
\cref{thm:univariate-gap-power} shows that the cleared value is still a product of
integer gap factors.  The computational table therefore confirms, rather than evades, the
one-direction obstruction.

For comparison, the positive transverse masks
\begin{equation}
 F_r(U,V)=(U-1)^{2r}+(V-1)^{2r}
 \label{eq:positive-transverse-mask}
\end{equation}
have exact contact $2r$, bidegree $(2r,2r)$, and support $4r+1$.  Their ordinary
coefficient $\ell^1$ norm is $2^{2r+1}$, while the shifted mask is simply
$X^{2r}+Y^{2r}$ and therefore has shifted norm $2$.  The first values are
\begin{center}
\small
\begin{tabular}{r r r r r}
\toprule
$r$ & contact & support & ordinary $\ell^1$ & shifted $\ell^1$ \\
\midrule
1&2&5&8&2\\
2&4&9&32&2\\
3&6&13&128&2\\
4&8&17&512&2\\
\bottomrule
\end{tabular}
\end{center}
They are nonzero by positivity and have an exceptionally small shifted norm.  At a shared
denominator prime, however, \cref{ex:positive-shared-mask} shows that unequal pole depths leave
the larger depth completely visible, while equal depths merely create one explicit two-term
residue test.  This is a compact calibration of the main tradeoff: nonvanishing and contact
are easy; repeated shared-prime denominator compression is the hard resource.

The cyclotomic-faced examples
\[
 (V+1)(U-1)^{2r}+(U+1)(V-1)^{2r}
\]
were checked symbolically for $1\le r\le4$.  Their contact is exactly $2r$, and their top
faces are $V+1$ and $U+1$ once $2r>1$.  They are among the smallest explicit positive examples in which a finite private prime can
be canceled by a genuine order condition on the opposite coordinate.  Stable support shows
that an asymptotic construction must combine such factors into tied multi-term faces at the
same structural primes, rather than assign fresh primes between coordinates.

\subsection{Reproducible source}

The following compact script reproduces the continued-fraction and uncertainty table.  It
uses high precision for reconnaissance; a proof-producing implementation would replace the
floating comparisons by directed logarithm intervals.

\begin{lstlisting}[style=pythonstyle,caption={Continued-fraction and uncertainty reconnaissance.}]
import mpmath as mp

mp.mp.dps = 100
log2, log3 = mp.log(2), mp.log(3)
theta = log3 / log2

# Continued fraction and convergents.
a = []
y = theta
for _ in range(18):
    z = int(mp.floor(y))
    a.append(z)
    y = 1 / (y - z)

p_m2, p_m1 = 0, 1
q_m2, q_m1 = 1, 0
conv = []
for z in a:
    p = z * p_m1 + p_m2
    q = z * q_m1 + q_m2
    delta = q * log3 - p * log2
    conv.append((p, q, delta))
    p_m2, p_m1 = p_m1, p
    q_m2, q_m1 = q_m1, q

for n, (p, q, delta) in enumerate(conv):
    if n == 0:
        ratio = None
    else:
        pp, qq, dd = conv[n - 1]
        H = max(p, q, pp, qq)
        ratio = 2 * H * max(abs(delta), abs(dd)) / log3
    print(n, p, q, delta, q * abs(delta), ratio)
\end{lstlisting}

For symbolic contact checks, one may expand $F(1+X,1+Y)$ and find the least total degree of a
nonzero monomial.  Toric masks must first be reduced modulo $I_{\boldsymbol\nu}$; otherwise an
identically zero lattice relation can masquerade as arbitrarily high contact.

\begin{statusbox}{\statuscomp{} for the symbolic polynomial identities; numerical
continued-fraction values are reconnaissance only.  No theorem in the report depends on the
printed decimals.}
\end{statusbox}

\section{Lean-oriented formalization plan}
\label{sec:lean-plan}

The new framework is unusually suitable for formalization because its core is finite algebra,
integer valuation bookkeeping, and continued-fraction determinant identities.  No Four
Exponentials input is needed to kernel-check the exact reductions.  A natural module split is
as follows.

\begin{longtable}{>{\raggedright\arraybackslash}p{0.25\textwidth}>{\raggedright\arraybackslash}p{0.67\textwidth}}
\toprule
Proposed module & Main statements \\
\midrule
\endhead
\texttt{RelationLattice} & Relation vectors, $\delta_\nu$, determinant identities
\eqref{eq:uncertainty-identities}, adjacent-convergent orientation, and the two-direction
uncertainty bound. \\
\texttt{MonomialToricKernel} & Finite-support monomial map, aggregation by exponent vector,
universal kernel equivalence, and invariance under unimodular basis changes. \\
\texttt{ContactMoments} & Shifted-polynomial contact, falling-factorial moments, ordinary
moment equivalence, degree ceiling, support ceiling, and the local $\ell^1$ estimate. \\
\texttt{UnivariateGapPower} & Divisibility of $S^dP(R/S)$ by $(R-S)^r$, including Laurent
normalization and the Thue--Morse specialization. \\
\texttt{BivariateGapIdeal} & Cleared evaluation formula \eqref{eq:cleared-bivariate},
membership in $(G_1,G_2)^r$, gcd-power divisibility, and the primewise tropical floor. \\
\texttt{PrivateLeadingFace} & The finite private-prime prototype: unique-minimum leading
face, unit-face survival, and exact first-layer cancellation. \\
\texttt{CyclotomicFiniteFace} & Homogenized cyclotomic evaluation, finite-field order
criterion, positivity of cross-faced masks, and finite-depth valuation conditions. \\
\texttt{StableSupport} & Equations \eqref{eq:structural-valuation}--
\eqref{eq:stable-reduced-support}, nonemptiness of $\mathcal P_\pm$, linear pole depths, and
adjacent valuation transport \eqref{eq:q-valuation-transport}--
\eqref{eq:p-valuation-transport}. \\
\texttt{SharedPoleNewton} & Exact term weights \eqref{eq:shared-newton-weight}, unique-minimum
equality, initial-form cancellation, unit-corner survival, and the positive-mask calibration.
\\
\texttt{ReducedGapTransport} & Lemma~\ref{lem:reduced-gap-transport}, structural-prime
exclusion, and the shared-gap mask corollary. \\
\texttt{ReducedDenominator} & Scalar criterion, compression factor, exact
shared-prime denominator law \eqref{eq:exact-shared-denominator}, and the finite implication
from a certified mask to contradiction. \\
\bottomrule
\end{longtable}

A low-risk theorem order is:
\begin{enumerate}
\item formalize contact, gap-ideal algebra, and toric normal forms over an arbitrary
  commutative domain;
\item specialize to $\Z$, rational evaluations, and prime valuations;
\item formalize continued fractions, support stabilization, and the determinant-one transport
  laws;
\item add the shared-pole Newton fan and exact reduced-denominator formula;
\item only then connect the finite certificate to real powers and effective logarithmic
  recurrence.
\end{enumerate}

The analytic interfaces can remain explicit hypotheses at first:
\begin{itemize}
\item $z_i=e^{\eta_i}$ and $|e^{\eta_i}-1|\le2|\eta_i|$ for $|\eta_i|\le1/4$;
\item effective recurrence $q_{n+1}\le Cq_n^\kappa$;
\item exact positivity or nonvanishing of the selected shared-face mask.
\end{itemize}
The first and third are routine real algebra/analysis once a concrete mask is given; the
second can be imported through a named linear-forms interface.  This separation keeps the
trust boundary visible.

A proof-producing search can output a compact certificate containing:
\begin{enumerate}
\item the exponent support and integer coefficients of $F$;
\item exact moment equations certifying contact and a toric normal form proving the mask is
  not an identically vanishing lattice relation;
\item a factorization, positivity, or Sturm certificate proving nonvanishing;
\item exact reduced fractions $z_n=R_n/S_n$, $z_{n+1}=R_{n+1}/S_{n+1}$ and the stable sets
  $\mathcal P_\pm$;
\item the cleared integer $J_F$ and its valuation at every prime in $\mathcal P_-$, including
  the successive shared-face initial-form cancellations;
\item the exact reduced denominator $b_F$ from \eqref{eq:exact-shared-denominator};
\item rational interval bounds verifying \eqref{eq:scalar-mask-criterion}.
\end{enumerate}
Every item except the final exponential estimate is finite integer arithmetic.  The final
estimate can use directed intervals for $\log2$ and $\log3$.

\begin{statusbox}{\statustarget{} only.  No new Lean files are claimed in this report.  The
listed statements are designed as independent additions rather than modifications of the
existing corpus.}
\end{statusbox}

\section{Public-status audit and literature placement}
\label{sec:public-audit}

The competitive framing in the prompt contains claims of very recent AI-assisted results.
Because none of those claims is a mathematical premise for the present report, the correct
standard is conservative: distinguish a public, inspectable result from an announcement, a
provisional solution, or a private rumor.

\subsection{What could be checked publicly on 22 August 2026}

A short explicit three-dimensional counterexample to the Jacobian conjecture was publicly
announced by Levent Alp\"oge with credit to Claude Fable, and independent expository accounts
verify the displayed map directly; the two-dimensional case remains open
\cite{TaoJacobian2026,GaoJacobian2026}.  This is a genuine public mathematical
development, but it has no known implication for the logarithmic period relation studied here.

Epoch AI's FrontierMath page for the twelve unresolved Hadamard orders below $2000$ reports
that Claude Fable 5.2 Pro submitted solutions for eleven orders and that a twelfth construction
was found but not yet fully written up; the page explicitly labels the status provisional
\cite{EpochHadamard2026}.  The corresponding Epoch page for elliptic curves of large rank still
listed the problem as unsolved and recorded rank at least $29$, not a publicly certified rank
$30$ construction, when accessed for this report \cite{EpochElliptic2026}.  Public posts and
news reports may be ahead of that page, but without a curve, independent points, and a
checkable rank certificate, the claimed rank-$30$ development is not used here.

No public paper, preprint, Lean repository, theorem statement, or inspectable proof describing
the rumored new \AEname\ breakthrough was located in the web audit.  Absence of a search result is
not evidence that no private work exists.  It means only that there is no public technical
input to import.  Accordingly, every mathematical claim in this report is proved here, marked
as an external theorem, or explicitly stated as a target.

\subsection{How the new method sits relative to known work}

The algebraic kernel in \cref{sec:toric} is part of the classical theory of lattice
and binomial ideals \cite{EisenbudSturmfels1996,Sturmfels1996}.  The new point is the use of
that kernel as a \emph{sanity filter} for continued-fraction near-unit masks: exact recurrence
relations among many convergents are toric identities and hence vanish on both the input and
output grids.

Prouhet--Tarry--Escott and Thue--Morse constructions are classical sources of moment
cancellation \cite{AlloucheShallit2003}.  The exact cleared-gap theorem
\cref{thm:univariate-gap-power} identifies why every such one-direction construction is
arithmetically inert here.  The bivariate gap ideal, stable structural-support theorem, shared-pole Newton law, exact
primewise denominator formula, and shared-face capture target are the genuinely
problem-specific additions.

The effective recurrence estimate \eqref{eq:effective-two-log} belongs to the theory of linear
forms in logarithms \cite{Baker1975,LaurentMignotteNesterenko1995}.  The reduced-gap estimate in
\cref{cor:nonadjacent-gcd-barrier} uses Bugeaud--Corvaja--Zannier
\cite{BugeaudCorvajaZannier2003}.  These theorems control the size of available near relations
and shared divisors; they do not address the bilinear logarithmic determinant directly.

The supplied unified report already explains the Four Exponentials barrier and develops
interpolation determinants, Kummer theory, compact orbits, natural boundaries, moment
methods, local rank, Smith profiles, and numerous other directions
\cite{UnifiedReport2026}.  This supplement deliberately does not repeat those surveys.
A literal source audit found no section organized around Prouhet or Thue--Morse masks,
toric/lattice ideals, augmentation-order gap ideals, stable structural support, or the
shared-pole Newton fan developed here.  The present report is
therefore an extension rather than a rearrangement of that manuscript.

\begin{statusbox}{\statusexternal{} for the dated public-status statements and cited
literature.  None of the news claims is used as a proof input.}
\end{statusbox}

\section{Conclusions and prioritized next steps}
\label{sec:conclusion}

No unconditional proof of the \AEname\ conjecture is claimed.  The progress is a sharp
reduction of a previously unseparated auxiliary-function architecture and, most importantly,
an exact identification of the denominator mechanism that remains after the natural
private-prime model fails.

\subsection{What has been proved in this report}

\begin{enumerate}
\item \textbf{Exact height of every exponent-lattice near-unit.}
  The reduced numerator and denominator are governed by the valuation norm $L_{M,N}$, so a
  good real relation still has exponential rational height.

\item \textbf{Toric collapse of many-convergent identities.}
  Universal monomial identities among any number of relation coordinates form a lattice ideal
  independent of the multiplicatively independent base pair.  Exact continued-fraction
  recurrences therefore vanish on both the input and output grids and cannot create a nonzero
  small integer.

\item \textbf{Exact one-direction no-go theorem.}
  Contact of order $r$ at $1$ forces a factor $(T-1)^r$; after clearing denominators, the mask
  value contains the $r$th power of the corresponding integer output gap.  This closes all
  one-direction Prouhet, Thue--Morse, and finite-difference masks at once.

\item \textbf{Multivariate augmentation-to-gap transfer.}
  Contact order $r$ at $(1,1)$ places every cleared bivariate value in $(G_1,G_2)^r$, with an
  exact primewise tropical floor and a scalar product-formula criterion.

\item \textbf{Stable structural support.}
  For adjacent convergents, every prime dividing $MN$ has an eventually fixed numerator or
  denominator sign.  Both reduced denominators are supported on the same nonempty finite set
  $\mathcal P_-$, and every pole depth grows linearly with the convergent denominator.  Thus
  fresh/private structural primes are not the asymptotic resource.

\item \textbf{Shared-pole Newton law.}
  At each $\ell\in\mathcal P_-$, every cleared monomial has the exact weight
  \eqref{eq:shared-newton-weight}.  A unique minimum survives; a tied face cancels one layer
  exactly when its explicit initial form vanishes.  Near-total compression requires this
  cancellation to repeat through almost the full depth at every stable denominator prime.

\item \textbf{Exact reduced-denominator endgame.}
  Formula \eqref{eq:exact-shared-denominator} expresses the denominator as the sum of the
  uncanceled shared pole depths.  The scalar criterion forces their total weighted residual to
  be only $O((d+e)\log q)$, while the visible pole mass is $\Theta(q(d+e))$.

\item \textbf{Quantitative no-go results.}
  Full visible-denominator clearing cannot become subunit; shared reduced-gap divisibility is
  fixed-size for adjacent directions and subheight for nonadjacent directions; unit-corner and
  generic unique-face masks leave full pole layers.
\end{enumerate}

\subsection{The surviving research program}

The highest-value next step is a structured search over shared Newton faces, not a broader
search over arbitrary high-contact coefficients.

\begin{enumerate}
\item \textbf{Prime-coupled shared-face search.}
  For the finite set $\mathcal P_-$, impose the exact linear congruences of
  \cref{app:milp} so that the cleared numerator captures all but $O(\log q)$ layers at every
  prime.  Optimize the rigorous scalar margin
  \[
   r\log\frac1{2\Delta}
   -\log b_F
   -\log\normone{F(1+X,1+Y)}.
  \]
  A positive interval-certified margin rules out the hypothetical candidate.

\item \textbf{Exploit determinant-one transport.}
  The depth vectors at adjacent indices satisfy
  \eqref{eq:q-valuation-transport}--\eqref{eq:p-valuation-transport}.  Search supports should
  be adapted to those exact affine relations, rather than treating the two pole slopes as
  independent tropical weights.

\item \textbf{Build or exclude deep initial-form cascades.}
  Cyclotomic factors provide finite residue identities, but a successful construction must
  combine them into multi-term faces that cancel at the same structural primes through many
  powers.  Conversely, a resultant or height theorem proving the lower bound
  \eqref{eq:shared-face-lower-target} would eliminate the whole mask architecture.

\item \textbf{Exact certificate search and Lean replay.}
  Contact equations, toric reduction, shared-prime congruences, valuations, and the reduced
  denominator are finite integer data.  Directed real intervals are needed only for the final
  strict inequality, making independent replay and kernel verification realistic.
\end{enumerate}

The conceptual lesson is now more precise than ``denominator cancellation is hard.''
Continued fractions provide excellent near-units and Prouhet masks provide excellent contact,
but the prime supports do not move: they stabilize.  The scarce resource is simultaneous,
deep cancellation on a fixed collection of shared structural-prime Newton faces.  That is the
remaining target exposed by this report.

\begin{statusbox}{The conjecture remains open.  For a fixed hypothetical candidate, one
certificate satisfying \cref{target:shared-face-capture} and
\eqref{eq:scalar-mask-criterion} would rule it out; a construction valid for every candidate
would prove the conjecture.  A uniform lower bound of the form
\eqref{eq:shared-face-lower-target} would instead rule out the entire exponent-lattice mask
architecture.}
\end{statusbox}

\appendix

\section{The general \texorpdfstring{$k$}{k}-variable gap-ideal theorem}
\label{app:k-gap}

The bivariate theorem has a direct finite-dimensional extension useful for formalization and
search.

\begin{theorem}[General augmentation-to-gap transfer]
\label{thm:k-gap-transfer}
Let $F\in\Z[T_1,\ldots,T_k]$ have degree at most $d_i$ in $T_i$.  Suppose
\[
 F(1+W_1,\ldots,1+W_k)=
 \sum_{\boldsymbol\alpha}\gamma_{\boldsymbol\alpha}W_1^{\alpha_1}\cdots W_k^{\alpha_k}
\]
has no nonzero term with $|\boldsymbol\alpha|<r$.  For rational points
$z_i=R_i/S_i$ with $R_i,S_i\in\Z$, $S_i>0$, put $G_i=R_i-S_i$ and
\[
 J_F=\left(\prod_{i=1}^kS_i^{d_i}\right)F(z_1,\ldots,z_k).
\]
Then $J_F\in\Z$ and
\begin{equation}
 J_F=\sum_{|\boldsymbol\alpha|\ge r}
 \gamma_{\boldsymbol\alpha}
 \prod_{i=1}^kG_i^{\alpha_i}S_i^{d_i-\alpha_i}
 \in(G_1,\ldots,G_k)^r.
 \label{eq:k-gap-expansion}
\end{equation}
Consequently
\begin{equation}
 \gcd(G_1,\ldots,G_k)^r\mid J_F.
 \label{eq:k-gap-gcd}
\end{equation}
For every prime $p$,
\begin{equation}
 v_p(J_F)\ge
 \min_{\gamma_{\boldsymbol\alpha}\ne0}
 \left(v_p(\gamma_{\boldsymbol\alpha})+
 \sum_{i=1}^k\bigl(\alpha_iv_p(G_i)+(d_i-\alpha_i)v_p(S_i)\bigr)\right).
 \label{eq:k-gap-tropical}
\end{equation}
\end{theorem}

\begin{proof}
Substitute $W_i=G_i/S_i$ into the shifted expansion and multiply by
$\prod S_i^{d_i}$.  Every exponent satisfies $0\le\alpha_i\le d_i$, so each term is integral.
The total gap degree is at least $r$, proving ideal membership and
\eqref{eq:k-gap-gcd}.  The valuation of a sum is at least the minimum valuation of its terms,
which gives \eqref{eq:k-gap-tropical}.
\end{proof}

This theorem separates two phenomena.  Contact controls total membership in a gap ideal;
primewise cancellation can improve the tropical minimum only through ties.  A proof search
must therefore record both the Newton support and the residue of every tied initial form.

\section{Weighted contact and tropical duality}
\label{app:weighted-contact}

Let $\boldsymbol\lambda=(\lambda_1,\ldots,\lambda_k)\in\R_{\ge0}^k$.  Define the weighted
contact of a shifted mask by
\begin{equation}
 \ord_{\boldsymbol\lambda}(F)=
 \min_{\gamma_{\boldsymbol\alpha}\ne0}
 \boldsymbol\lambda\cdot\boldsymbol\alpha.
 \label{eq:weighted-contact}
\end{equation}
If $|z_i-1|\le e^{-t\lambda_i}$, then
\begin{equation}
 |F(z_1,\ldots,z_k)|
 \le\normone{F(1+\mathbf W)}e^{-t\ord_{\boldsymbol\lambda}(F)}.
 \label{eq:weighted-archimedean}
\end{equation}
At a prime $p\nmid S_1\cdots S_k$, taking
$\lambda_i=v_p(G_i)$ makes the same Newton support produce the lower bound
\begin{equation}
 v_p(J_F)\ge
 \min_{\gamma_{\boldsymbol\alpha}\ne0}
 \left(v_p(\gamma_{\boldsymbol\alpha})+
 \boldsymbol\lambda\cdot\boldsymbol\alpha\right).
 \label{eq:weighted-padic}
\end{equation}
Thus the archimedean gain and the nonarchimedean floor are two evaluations of the same Newton
polyhedron.  Automatic termwise divisibility cannot create a contradiction by itself; the
only surplus can come from different faces being active at different places, or from residue
cancellation on tied faces.  In the present setting, the leading-face theorem identifies the
finite places at which such a change of active face is compulsory.

\section{Integer-programming formulation of the mask search}
\label{app:milp}

Fix a finite exponent support
$E=\{(a_\ell,b_\ell):1\le\ell\le L\}\subset\Z_{\ge0}^2$ and unknown coefficients
$c_\ell\in\Z$.  Contact at least $r$ is the homogeneous integer system
\begin{equation}
 \sum_{\ell=1}^L c_\ell
 \binom{a_\ell}{i}\binom{b_\ell}{j}=0
 \qquad(i+j<r).
 \label{eq:milp-contact}
\end{equation}
The objective $\sum|c_\ell|$ is linearized with auxiliary variables.  A nonzero normalization
may fix one primitive coefficient or impose a signed face sum.  Toric reduction should be
performed before optimization.

For a fixed adjacent candidate pair, shared-face capture is also linear in the unknown
coefficients.  Let $d=\max a_\ell$, $e=\max b_\ell$ and write
\begin{equation}
 J_F=\sum_{\ell=1}^L c_\ell
 R_1^{a_\ell}R_2^{b_\ell}
 S_1^{d-a_\ell}S_2^{e-b_\ell}.
 \label{eq:milp-cleared-linear-form}
\end{equation}
For every $p\in\mathcal P_-$ choose an allowed residual depth $B_p$ and impose
\begin{equation}
 J_F\equiv0
 \pmod{p^{d s_{p,1}+e s_{p,2}-B_p}}.
 \label{eq:milp-shared-congruence}
\end{equation}
This is a linear congruence in the $c_\ell$; Chinese remaindering combines the finitely many
structural primes.  Taking $B_p=O(\log q)$ encodes the near-total capture forced by
\eqref{eq:shared-residual-polyheight}.  Positivity can be enforced by a cyclotomic or
sum-of-squares ansatz, or checked afterward by an exact Sturm procedure.

A practical search should not maximize contact alone.  Its exact target is the scalar margin
\begin{equation}
 \mathcal M(F)=
 r\log\frac1{2\Delta}
 -\log b_F
 -\log\normone{F(1+X,1+Y)}.
 \label{eq:search-margin}
\end{equation}
A positive rigorously bounded margin is already a proof certificate.  Negative margins reveal
which prime layers dominate the failure and can guide the next face support.

\section{Notation crosswalk}
\label{app:notation}

\begin{center}
\small
\begin{tabularx}{\textwidth}{@{}lX@{}}
\toprule
Symbol & Meaning \\
\midrule
$x$ & Hypothetical nonintegral exponent. \\
$M,N$ & Integer outputs $2^x,3^x$. \\
$\thetaAE$ & $\log3/\log2$. \\
$\nu=(a,b)$ & Exponent-lattice relation vector. \\
$\delta_\nu$ & $a\log2+b\log3$. \\
$\rho_\nu,\sigma_\nu$ & Input/output near-units $2^a3^b$ and $M^aN^b$. \\
$R_\nu/S_\nu$ & Reduced form of $\sigma_\nu$. \\
$G_\nu$ & Gap $R_\nu-S_\nu$. \\
$u_{\ell,k}$ & Structural valuation $q_kn_\ell-p_km_\ell$ of the $k$th near-unit. \\
$\mathcal P_\pm$ & Stable numerator/denominator prime supports. \\
$s_{\ell,k}$ & Stable denominator depth $-u_{\ell,k}$ for $\ell\in\mathcal P_-$. \\
$I_{\boldsymbol\nu}$ & Universal toric/lattice kernel of the monomial coordinate map. \\
$\ord_{\mathbf1}(F)$ & Contact order of a mask at $(1,1)$. \\
$J_F$ & Fully cleared integer mask value. \\
$\rho_{\ell,k}(F)$ & Uncanceled shared pole depth in the reduced denominator. \\
$b_F$ & Actual reduced denominator of the rational mask value. \\
$K_F$ & Compression factor from visible to reduced denominator. \\
$\mathfrak c$ & Normalized denominator-compression index. \\
$\Phi_m^{\mathrm{hom}}$ & Homogenized cyclotomic polynomial. \\
\bottomrule
\end{tabularx}
\end{center}

\begin{thebibliography}{99}

\bibitem{AE1944}
L.~Alaoglu and P.~Erd\H{o}s,
\newblock \emph{On highly composite and similar numbers},
\newblock Transactions of the American Mathematical Society \textbf{56} (1944), 448--469.

\bibitem{Baker1975}
A.~Baker,
\newblock \emph{Transcendental Number Theory},
\newblock Cambridge University Press, 1975.

\bibitem{LaurentMignotteNesterenko1995}
M.~Laurent, M.~Mignotte, and Yu.~Nesterenko,
\newblock \emph{Formes lin\'eaires en deux logarithmes et d\'eterminants d'interpolation},
\newblock Journal of Number Theory \textbf{55} (1995), 285--321.

\bibitem{BugeaudCorvajaZannier2003}
Y.~Bugeaud, P.~Corvaja, and U.~Zannier,
\newblock \emph{An upper bound for the G.C.D. of $a^n-1$ and $b^n-1$},
\newblock Mathematische Zeitschrift \textbf{243} (2003), 79--84.
\newblock DOI: \href{https://doi.org/10.1007/s00209-002-0449-z}{10.1007/s00209-002-0449-z}.

\bibitem{EisenbudSturmfels1996}
D.~Eisenbud and B.~Sturmfels,
\newblock \emph{Binomial ideals},
\newblock Duke Mathematical Journal \textbf{84} (1996), 1--45.
\newblock DOI: \href{https://doi.org/10.1215/S0012-7094-96-08401-X}{10.1215/S0012-7094-96-08401-X}.

\bibitem{Sturmfels1996}
B.~Sturmfels,
\newblock \emph{Gr\"obner Bases and Convex Polytopes},
\newblock University Lecture Series 8, American Mathematical Society, 1996.

\bibitem{AlloucheShallit2003}
J.-P.~Allouche and J.~Shallit,
\newblock \emph{Automatic Sequences: Theory, Applications, Generalizations},
\newblock Cambridge University Press, 2003.

\bibitem{Khinchin1964}
A.~Ya.~Khinchin,
\newblock \emph{Continued Fractions},
\newblock University of Chicago Press, 1964.

\bibitem{TaoJacobian2026}
T.~Tao,
\newblock \emph{A digestion of the Jacobian conjecture counterexample},
\newblock 21 July 2026,
\newblock \url{https://terrytao.wordpress.com/2026/07/21/a-digestion-of-the-jacobian-conjecture-counterexample/}.

\bibitem{GaoJacobian2026}
S.~Gao,
\newblock \emph{Counterexamples to the Jacobian conjecture in dimensions greater than two},
\newblock arXiv:2608.00222, 2026.

\bibitem{EpochHadamard2026}
Epoch AI,
\newblock \emph{Hadamard Matrices of Orders $n<2000$}, FrontierMath open-problem page,
\newblock accessed 22 August 2026,
\newblock \url{https://epoch.ai/frontiermath/open-problems/hadamard-matrices}.

\bibitem{EpochElliptic2026}
Epoch AI,
\newblock \emph{Elliptic Curves over $\mathbb Q$ of Large Rank}, FrontierMath open-problem page,
\newblock accessed 22 August 2026,
\newblock \url{https://epoch.ai/frontiermath/open-problems/elliptic-curve-rank}.

\bibitem{UnifiedReport2026}
The ProveIt project,
\newblock \emph{The Alaoglu--Erd\H{o}s Integer-Exponent Conjecture: Unified Report},
\newblock manuscript supplied with the present task, 22 August 2026.

\end{thebibliography}

\end{document}
